\documentclass[aspectratio=169,10pt]{beamer}
% Preamble for the Computational Intelligence slides (Beamer, Metropolis)
\usetheme[progressbar=frametitle,numbering=fraction,background=light,block=fill,sectionpage=none]{metropolis}
\usepackage[T1]{fontenc}
\usepackage{FiraSans}
\usepackage[scaled=.85]{FiraMono}
\usepackage{amsmath,amssymb}
\usefonttheme[onlymath]{serif}
\usepackage{booktabs,array,tabularx}
\usepackage{tikz,pgfplots}
\usepackage{xcolor}
\usepackage{listings}
\usepackage{adjustbox}
\usepackage{fontawesome5}
\hyphenpenalty=10000\exhyphenpenalty=10000
\newenvironment{sfig}[1][0.80]{\begin{adjustbox}{max width=\linewidth,max totalheight=#1\textheight,center}}{\end{adjustbox}}
\usetikzlibrary{arrows.meta,positioning,shapes.geometric,shapes.misc,fit,calc,backgrounds,decorations.pathreplacing,patterns,matrix,chains}
\pgfplotsset{compat=1.17}

\definecolor{tdInk}{HTML}{1D2330}
\definecolor{tdMuted}{HTML}{5D6678}
\definecolor{tdLine}{HTML}{D9D5CA}
\definecolor{tdPaper}{HTML}{FBFAF7}
\definecolor{tdBlue}{HTML}{2457D6}
\definecolor{tdGreen}{HTML}{17855A}
\definecolor{tdAmber}{HTML}{B7791F}
\definecolor{tdRed}{HTML}{C8362F}
\definecolor{tdPurple}{HTML}{7A4FD1}
\setbeamercolor{normal text}{fg=tdInk,bg=tdPaper}
\setbeamercolor{alerted text}{fg=tdRed}
\setbeamercolor{example text}{fg=tdGreen}
\setbeamercolor{frametitle}{bg=tdInk,fg=white}
\setbeamercolor{progress bar}{fg=tdBlue,bg=tdLine}
\setbeamercolor{title separator}{fg=tdBlue}
\setbeamercolor{block title}{bg=tdBlue!12,fg=tdInk}
\setbeamercolor{block body}{bg=tdBlue!5,fg=tdInk}
\setbeamercolor{block title alerted}{bg=tdRed!12,fg=tdRed!80!black}
\setbeamercolor{block body alerted}{bg=tdRed!5,fg=tdInk}
\setbeamercolor{block title example}{bg=tdGreen!14,fg=tdGreen!70!black}
\setbeamercolor{block body example}{bg=tdGreen!6,fg=tdInk}
\setbeamerfont{frametitle}{size=\large,series=\bfseries}
\makeatletter\setlength{\metropolis@frametitle@padding}{1.15ex}\makeatother
\setbeamertemplate{navigation symbols}{}
\setbeamerfont{block title}{size=\small}

% ---- code listings (Python 3.14) ----
\lstdefinestyle{py}{language=Python,
  basicstyle=\ttfamily\fontsize{8.4}{10.1}\selectfont\color{tdInk},
  keywordstyle=\color{tdBlue}\bfseries, commentstyle=\color{tdMuted}\itshape, stringstyle=\color{tdGreen},
  morekeywords={match,case,type,self,True,False,None},
  emph={print,len,range,sorted,min,max,sum,enumerate,zip,heappush,heappop,cache},emphstyle=\color{tdPurple},
  showstringspaces=false, columns=fullflexible, keepspaces=true, frame=single, rulecolor=\color{tdLine},
  backgroundcolor=\color{white}, xleftmargin=2pt, xrightmargin=2pt, aboveskip=2pt, belowskip=0pt, framesep=3pt}
\lstdefinestyle{out}{basicstyle=\ttfamily\fontsize{8}{9.6}\selectfont\color{tdInk}, columns=fullflexible, keepspaces=true,
  frame=single, rulecolor=\color{tdGreen!60}, backgroundcolor=\color{tdGreen!6}, xleftmargin=2pt, xrightmargin=2pt, aboveskip=2pt, belowskip=0pt, framesep=3pt}
% \pycode[size]{snippet} shows snippets/<snippet>.py ; \pyout[size]{snippet} shows snippets/<snippet>.out
\newcommand{\pycode}[2][\relax]{{#1\lstinputlisting[style=py]{snippets/#2.py}}}
\newcommand{\pyout}[2][\relax]{{#1\lstinputlisting[style=out]{snippets/#2.out}}}

% ---- tikz styles ----
\tikzset{>={Stealth[length=2.2mm]},
  card/.style={rounded corners=4pt,draw=#1!70!black,fill=#1!10,align=center,inner sep=5pt,font=\small}, card/.default=tdBlue,
  pill/.style={rounded corners=7pt,draw=none,fill=#1,text=white,inner sep=3pt,font=\scriptsize\bfseries,align=center}, pill/.default=tdBlue,
  arr/.style={->,thick,draw=tdMuted},
  cell/.style={draw=tdInk!70,minimum width=7.5mm,minimum height=7.5mm,inner sep=0pt,font=\small\ttfamily,fill=white},
  tn/.style={circle,draw=tdInk!80,fill=white,minimum size=7mm,inner sep=0pt,font=\small\bfseries}}
\newcommand{\key}[1]{\textbf{\color{tdBlue}#1}}
\newcommand{\bad}[1]{\textbf{\color{tdRed}#1}}
\newcommand{\good}[1]{\textbf{\color{tdGreen!70!black}#1}}
\newcommand{\icsfooter}[1]{\setbeamertemplate{footline}{\begin{tikzpicture}[remember picture,overlay]\node[anchor=south east,font=\tiny,text=tdMuted,inner sep=0pt] at ([xshift=-7mm,yshift=1.6mm]current page.south east) {Introduction to Computing Science \,|\, #1 \,|\, BUV\hspace{1.2em}\insertframenumber/\inserttotalframenumber};\end{tikzpicture}}}
\newcommand{\icstitle}[2]{\title{#1}\subtitle{#2}\author{Dr.\ Quang-Vinh Dang\\ {\small Senior Lecturer (Assoc.\ Prof)}}\institute{British University Vietnam}\date{Introduction to Computing Science \,|\, BSc Data Science and AI \,|\, Python 3.14}}
\usepackage{multicol}
\pgfplotsset{every axis/.append style={tick label style={font=\scriptsize},label style={font=\scriptsize},legend style={font=\scriptsize,draw=none,fill=none}}}
\tikzset{soft/.style={rounded corners=4pt,draw=tdLine,fill=white,align=center,inner sep=4pt,font=\small},
  lbl/.style={font=\scriptsize,text=tdMuted,align=center}}
% source footnote at the bottom of a slide
\newcommand{\src}[1]{\par\vfill{\fontsize{5.5}{6.4}\selectfont\color{tdMuted}#1}\vspace{1mm}}
% discussion prompt bar
\newcommand{\discuss}[1]{\par\vfill\noindent\begin{tikzpicture}\node[rounded corners=4pt,draw=tdAmber!80!black,fill=tdAmber!12,inner sep=5pt,text width=\linewidth-12pt,align=left,font=\small]{\textbf{\color{tdAmber!80!black}Discuss.}\ #1};\end{tikzpicture}\par}
% philosophy slide: \philslide{title}{source}{framing}{A label}{A}{B label}{B}{voices}{discussion}; sketch k = \psk<k>
\newcounter{pscount}
\newcommand{\defsk}[2]{\expandafter\gdef\csname psk#1\endcsname{#2}}
\tikzset{sk/.style={rounded corners=3pt,draw=#1!70!black,fill=#1!10,inner sep=2.5pt,font=\scriptsize,align=center,text width=2.3cm},sk/.default=tdBlue}
\newcommand{\philslide}[9]{\stepcounter{pscount}\begin{frame}{#1}
\begin{columns}[T]
\begin{column}{0.5\textwidth}
{\footnotesize #3}\par\vspace{1mm}
\begin{sfig}[0.18]\begin{tikzpicture}\csname psk\arabic{pscount}\endcsname\end{tikzpicture}\end{sfig}
\end{column}
\begin{column}{0.47\textwidth}
\begin{tikzpicture}\node[rounded corners=4pt,draw=tdGreen!70!black,fill=tdGreen!9,inner sep=4pt,text width=\linewidth-10pt,align=left,font=\footnotesize]{\textbf{\color{tdGreen!60!black}#4}\\[1pt] #5};\end{tikzpicture}\par\vspace{1.5mm}
\begin{tikzpicture}\node[rounded corners=4pt,draw=tdPurple!70!black,fill=tdPurple!8,inner sep=4pt,text width=\linewidth-10pt,align=left,font=\footnotesize]{\textbf{\color{tdPurple!70!black}#6}\\[1pt] #7};\end{tikzpicture}\par\vspace{1mm}
{\scriptsize\color{tdMuted}\textbf{Key voices:} #8}
\end{column}
\end{columns}
\discuss{#9}
\src{#2}
\end{frame}}
\newcommand{\qcard}[3]{\begin{tikzpicture}\node[rounded corners=4pt,draw=#1!70!black,fill=#1!8,inner sep=4pt,text width=#2,align=left,font=\footnotesize]{#3};\end{tikzpicture}}
\newcommand{\illus}{\textit{illustrative}}

\lstdefinestyle{pys}{style=py,basicstyle=\ttfamily\fontsize{7.3}{8.5}\selectfont\color{tdInk}}
\lstdefinestyle{outs}{style=out,basicstyle=\ttfamily\fontsize{7.3}{8.8}\selectfont\color{tdInk}}
\newcommand{\pycodes}[1]{\lstinputlisting[style=pys]{snippets/#1.py}}
\newcommand{\pyouts}[1]{\lstinputlisting[style=outs]{snippets/#1.out}}

\icsfooter{Computational Intelligence: What is AI?}
\icstitle{What is AI?}{Intelligence, minds and machines}

\newcommand{\agn}[4]{\node[card=#3,minimum width=2.45cm,minimum height=1.9cm,font=\footnotesize] (n#1) at ({mod(#1-1,5)*2.9},{-floor((#1-1)/5)*2.7}) {\textbf{#1. #2}\\[3pt]\shortstack{\scriptsize\color{tdMuted}#4}};}
\begin{document}
\begin{frame}[plain]\titlepage\end{frame}

\begin{frame}{Where this lecture goes}
\begin{sfig}[0.78]
\begin{tikzpicture}[x=1cm,y=1cm]
  \agn{1}{Definitions}{tdBlue}{what is intelligence?\\what is AI?}
  \agn{2}{History}{tdGreen}{1956 to\\agents}
  \agn{3}{Turing test}{tdAmber}{imitation game,\\new results}
  \agn{4}{Chinese Room}{tdRed}{syntax vs\\meaning}
  \agn{5}{Tests}{tdPurple}{ARC-AGI,\\HLE}
  \agn{6}{Big questions}{tdBlue}{mind, ethics,\\risk}
  \agn{7}{Film and TV}{tdGreen}{1927 to 2026}
  \agn{8}{AI now}{tdAmber}{late 2026}
  \agn{9}{You}{tdRed}{debates,\\exercise}
  \foreach \a/\b in {1/2,2/3,3/4,4/5,6/7,7/8,8/9}{\draw[arr] (n\a) -- (n\b);}
  \draw[arr] (n5.south) -- ++(0,-0.55) -| (n6.north);
\end{tikzpicture}
\end{sfig}
\vspace{1mm}
{\small Intuition first, no prerequisites. Three short Python 3.14 demos (ELIZA, a Chinese Room, a ``parrot'') with their real output.}
\end{frame}

\begin{frame}{After this lecture you can}
\begin{columns}[T]
\begin{column}{0.54\textwidth}
\begin{itemize}\setlength\itemsep{4pt}\small
\item name the \textbf{four classic views} of AI and say which one modern AI mostly follows;
\item compare three definitions of intelligence (\textbf{goals}, \textbf{skill-acquisition efficiency}, \textbf{cognitive breadth});
\item explain the \textbf{Turing test}, why \textbf{ELIZA} fooled people, and what the \textbf{2025--26 results} do and do not show;
\item state \textbf{Searle's Chinese Room} and the standard replies;
\item discuss whether \textbf{LLMs understand}, with evidence on both sides;
\item pose \textbf{two dozen philosophical questions} about minds and machines, and defend a position;
\item connect these questions to \textbf{films and series} from 1927 to 2026.
\end{itemize}
\end{column}
\begin{column}{0.42\textwidth}
\begin{block}{How to read the slides}
{\small
\begin{itemize}\setlength\itemsep{2pt}
\item Amber boxes are \textbf{discussion prompts}.
\item Small grey text at the bottom is the \textbf{source}.
\item Anything about 2025--2026 comes from a source I found; where I could not verify, the slide says so.
\item Philosophy slides are \textbf{balanced on purpose}: pick a side and argue.
\end{itemize}}
\end{block}
\end{column}
\end{columns}
\end{frame}

% ---------- 1. Definitions ----------
\begin{frame}{Warm-up: which of these is ``AI''?}
\begin{columns}[T]
\begin{column}{0.55\textwidth}
\small
Vote with your hand for each item. There is no right answer yet.
\begin{enumerate}\setlength\itemsep{3pt}
\item A thermostat that keeps a room at 22\,$^\circ$C.
\item A spam filter that learns from your ``report spam'' clicks.
\item A chess program that beats every human.
\item A calculator.
\item A chatbot that writes essays.
\item A robot vacuum that maps your flat.
\item A search engine ranking pages.
\end{enumerate}
\end{column}
\begin{column}{0.42\textwidth}
\qcard{tdBlue}{\linewidth-10pt}{\textbf{Notice what you used to decide}\\[2pt]
Did it \emph{learn}? Does it \emph{act} sensibly? Is it \emph{hard for humans}? Does it \emph{seem} to understand?\\[2pt]
These are different definitions of AI. The rest of this part makes them precise.}
\vspace{2mm}\par
\qcard{tdAmber}{\linewidth-10pt}{\textbf{The ``AI effect''}\\[2pt]
Tasks that once counted as AI (chess, route planning, OCR) are often renamed ``just software'' once they work. The target keeps moving.}
\end{column}
\end{columns}
\discuss{Where did your class disagree most? What does that say about the word ``intelligence''?}
\end{frame}

\begin{frame}{What is intelligence? Even experts do not agree}
\begin{columns}[T]
\begin{column}{0.49\textwidth}
\qcard{tdBlue}{\linewidth-10pt}{\textbf{Wechsler (1944), psychology}\\ ``the aggregate or global capacity of the individual to act purposefully, to think rationally and to deal effectively with his environment.''}\par\vspace{2mm}
\qcard{tdGreen}{\linewidth-10pt}{\textbf{Gottfredson (1997), psychology}\\ ``a very general mental capability that, among other things, involves the ability to reason, plan, solve problems, think abstractly, comprehend complex ideas, learn quickly and learn from experience.''}
\end{column}
\begin{column}{0.49\textwidth}
\qcard{tdAmber}{\linewidth-10pt}{\textbf{Legg and Hutter (2007), AI}\\ ``Intelligence measures an agent's ability to achieve goals in a wide range of environments.''}\par\vspace{2mm}
\qcard{tdPurple}{\linewidth-10pt}{\textbf{Chollet (2019), AI}\\ ``The intelligence of a system is a measure of its skill-acquisition efficiency over a scope of tasks, with respect to priors, experience, and generalization difficulty.''}
\end{column}
\end{columns}
\vspace{1mm}
{\small Common threads: \key{goals}, \key{adapting}, \key{learning}, \key{breadth}. Legg and Hutter collected \textbf{over 70} published definitions before proposing their own.}
\src{Wechsler, \emph{The Measurement of Adult Intelligence} (1944); Gottfredson, Intelligence 24(1), 1997; Legg and Hutter, ``Universal Intelligence: A Definition of Machine Intelligence'', Minds and Machines 17, 2007; Chollet, ``On the Measure of Intelligence'', arXiv:1911.01547, 2019.}
\end{frame}

\begin{frame}{Four views of AI (Russell and Norvig)}
\begin{sfig}[0.62]
\begin{tikzpicture}
  \node[lbl,font=\footnotesize\bfseries] at (2.6,2.3) {Human-like};
  \node[lbl,font=\footnotesize\bfseries] at (8.1,2.3) {Rational (ideal)};
  \node[lbl,font=\footnotesize\bfseries,rotate=90] at (-0.8,0.6) {Thinking};
  \node[lbl,font=\footnotesize\bfseries,rotate=90] at (-0.8,-2.0) {Acting};
  \node[card=tdBlue,text width=4.9cm,minimum height=2.1cm] at (2.6,0.6) {\textbf{Thinking humanly}\\ model how people think\\ {\scriptsize cognitive science, brain models}};
  \node[card=tdGreen,text width=4.9cm,minimum height=2.1cm] at (8.1,0.6) {\textbf{Thinking rationally}\\ ``laws of thought'', logic\\ {\scriptsize Aristotle's syllogisms, theorem provers}};
  \node[card=tdAmber,text width=4.9cm,minimum height=2.1cm] at (2.6,-2.0) {\textbf{Acting humanly}\\ behave so people cannot tell\\ {\scriptsize the Turing test}};
  \node[card=tdRed,text width=4.9cm,minimum height=2.1cm] at (8.1,-2.0) {\textbf{Acting rationally}\\ do the best thing given what is known\\ {\scriptsize the rational agent}};
  \node[pill=tdRed,anchor=north] at (8.1,-3.25) {the view most of modern AI follows};
\end{tikzpicture}
\end{sfig}
\src{Russell and Norvig, \emph{Artificial Intelligence: A Modern Approach}, 4th ed., Pearson, 2020, ch.\ 1.}
\end{frame}

\begin{frame}{Why ``acting rationally'' won}
\begin{columns}[T]
\begin{column}{0.5\textwidth}
\small
\begin{itemize}\setlength\itemsep{4pt}
\item \textbf{Measurable.} Rationality is defined mathematically: maximise expected performance.
\item \textbf{General.} Covers logic, learning, perception, control.
\item \textbf{Not tied to humans.} People are not perfectly rational, and we may not want machines to copy our mistakes.
\end{itemize}
\vspace{1mm}
Analogy used by Russell and Norvig: aeronautical engineers did not succeed by building machines that fly exactly like pigeons.
\end{column}
\begin{column}{0.47\textwidth}
\begin{sfig}[0.55]
\begin{tikzpicture}
  \node[card=tdBlue,minimum width=2.6cm,minimum height=1.7cm] (ag) at (0,0) {\textbf{Agent}\\\shortstack{\scriptsize chooses the action\\ that maximises\\ expected performance}};
  \node[card=tdGreen,minimum width=2.6cm,minimum height=1.7cm] (en) at (5,0) {\textbf{Environment}\\\shortstack{\scriptsize world, users,\\ other agents}};
  \draw[arr] (en.north west) to[bend right=25] node[above,lbl] {percepts (sensors)} (ag.north east);
  \draw[arr] (ag.south east) to[bend right=25] node[below,lbl] {actions (actuators)} (en.south west);
  \node[lbl] at (2.5,-1.7) {performance measure: what counts as ``good''};
\end{tikzpicture}
\end{sfig}
\end{column}
\end{columns}
\discuss{A self-driving car is a rational agent. Who chooses its performance measure, and what should it include?}
\end{frame}

\begin{frame}{Legg and Hutter: intelligence as goal achievement}
\begin{columns}[T]
\begin{column}{0.52\textwidth}
\small
Idea: an agent is more intelligent if it does well in \textbf{more kinds of worlds}, giving more weight to \textbf{simple} worlds.
\begin{equation*}
\Upsilon(\pi)=\sum_{\mu\in E} 2^{-K(\mu)}\,V^{\pi}_{\mu}
\end{equation*}
{\footnotesize $\pi$: the agent; $\mu$: an environment; $K(\mu)$: its description length (Kolmogorov complexity); $V^{\pi}_{\mu}$: expected reward of $\pi$ in $\mu$.}
\vspace{2mm}

\textbf{Good:} one number, no mention of humans or brains.\\
\textbf{Problem:} $K$ is uncomputable, so it is an \emph{ideal}, not a test you can run.
\end{column}
\begin{column}{0.45\textwidth}
\begin{sfig}[0.6]
\begin{tikzpicture}
  \draw[->] (0,0) -- (5.2,0) node[right,font=\scriptsize] {environments};
  \draw[->] (0,0) -- (0,3.2) node[above,font=\scriptsize] {score};
  \foreach \x/\h in {0.4/2.5,1.1/2.1,1.8/1.2,2.5/0.9,3.2/0.5,3.9/0.4,4.6/0.2}{\fill[tdBlue!35] (\x,0) rectangle (\x+0.5,\h);}
  \foreach \x/\h in {0.4/2.9,1.1/2.8,1.8/2.6,2.5/2.5,3.2/2.2,3.9/2.0,4.6/1.8}{\draw[tdGreen,very thick] (\x,\h) -- (\x+0.5,\h);}
  \node[font=\scriptsize,text=tdGreen!60!black,anchor=west] at (2.6,3.1) {broad agent};
  \node[font=\scriptsize,text=tdBlue,anchor=west] at (2.6,2.7) {narrow agent (bars)};
  \node[lbl] at (2.6,-0.5) {simple worlds \hspace{2.4cm} complex worlds};
\end{tikzpicture}
\end{sfig}
{\scriptsize\illus}
\end{column}
\end{columns}
\src{Legg and Hutter, ``Universal Intelligence: A Definition of Machine Intelligence'', Minds and Machines 17(4), 2007.}
\end{frame}

\begin{frame}{Chollet: intelligence is efficient skill acquisition}
\begin{columns}[T]
\begin{column}{0.5\textwidth}
\small
\textbf{Skill is not intelligence.} A system can buy skill with \emph{priors} (built-in knowledge) or \emph{experience} (huge training data).

Chollet's proposal: intelligence = how much new skill you gain per unit of priors and experience, over tasks that are \textbf{new} to you (generalisation).

\vspace{1mm}
Consequence: beating humans at chess after millions of games shows \textbf{skill}, but says little about the \textbf{efficiency} with which the skill was acquired.
\end{column}
\begin{column}{0.47\textwidth}
\begin{sfig}[0.62]
\begin{tikzpicture}
  \node[lbl,font=\footnotesize\bfseries] at (1.0,3.4) {Chess engine};
  \node[lbl,font=\footnotesize\bfseries] at (4.2,3.4) {Child};
  \fill[tdRed!30] (0.3,0) rectangle (1.7,2.8); \node[lbl] at (1.0,1.4) {experience:\\ millions of\\ games};
  \fill[tdGreen!30] (3.5,0) rectangle (4.9,0.8); \node[lbl,font=\tiny] at (4.2,0.4) {few\\ examples};
  \node[pill=tdRed] at (1.0,-0.6) {high skill, low efficiency};
  \node[pill=tdGreen] at (4.2,-0.6) {lower skill, high efficiency};
  \draw[tdMuted,dashed] (-0.2,0) -- (5.3,0);
\end{tikzpicture}
\end{sfig}
{\scriptsize\illus: bar height = experience consumed.}
\end{column}
\end{columns}
\src{Chollet, ``On the Measure of Intelligence'', arXiv:1911.01547, 2019. ARC-AGI (later in this lecture) was designed to measure this.}
\end{frame}

\begin{frame}{AGI: three definitions used today}
\begin{sfig}[0.62]
\begin{tikzpicture}
  \node[card=tdBlue,text width=5.2cm,minimum height=2.6cm,anchor=north west] at (0,0) {\textbf{Economic}\\[2pt] \footnotesize ``highly autonomous systems that outperform humans at most economically valuable work''\\[3pt]{\scriptsize OpenAI Charter (2018)}};
  \node[card=tdGreen,text width=5.2cm,minimum height=2.6cm,anchor=north west] at (5.6,0) {\textbf{Levels of AGI}\\[2pt] \footnotesize two axes: \emph{performance} (emerging, competent, expert, virtuoso, superhuman) and \emph{generality} (narrow or general)\\[3pt]{\scriptsize Morris et al., Google DeepMind (2023)}};
  \node[card=tdAmber,text width=5.2cm,minimum height=2.6cm,anchor=north west] at (11.2,0) {\textbf{Cognitive profile}\\[2pt] \footnotesize AGI = matching the cognitive versatility and proficiency of a well-educated adult, scored over ten domains (CHC theory)\\[3pt]{\scriptsize Hendrycks et al. (2025)}};
\end{tikzpicture}
\end{sfig}
\vspace{1mm}
\begin{columns}[T]
\begin{column}{0.58\textwidth}
\small
The 2025 paper reports an ``AGI score'': \textbf{GPT-4 (2023) 27\%}, \textbf{GPT-5 (2025) 57\%}, with long-term memory storage scoring 0\% for GPT-5. A number is only as good as the definition behind it.
\end{column}
\begin{column}{0.38\textwidth}
\qcard{tdRed}{\linewidth-10pt}{\scriptsize Note: no definition is agreed. Chollet's (skill-acquisition efficiency) is a fourth, and stricter, option.}
\end{column}
\end{columns}
\src{OpenAI Charter, openai.com/charter; Morris et al., ``Levels of AGI'', arXiv:2311.02462; Hendrycks et al., ``A Definition of AGI'', arXiv:2510.18212 (Oct 2025, rev.\ Dec 2025), Table 1.}
\end{frame}

\begin{frame}{Narrow, general, super: a popular ladder}
\begin{sfig}[0.6]
\begin{tikzpicture}
  \foreach \i/\t/\d/\c in {0/Narrow AI (ANI)/{one task or domain: spam filter, chess, translation}/tdBlue,1/General AI (AGI)/{broad human-level competence, learns new tasks}/tdGreen,2/Superintelligence (ASI)/{far beyond the best humans in most domains}/tdPurple}{
    \node[card=\c,minimum width={4.2cm+\i*2.6cm},minimum height=1.3cm,anchor=south west,font=\small] at (0,{\i*1.55}) {\textbf{\t}\\{\scriptsize\d}};
  }
  \draw[->,very thick,tdMuted] (-0.5,0) -- (-0.5,4.5) node[above,lbl] {capability};
  \node[lbl,anchor=west] at (9.2,0.6) {today's systems?};
  \node[lbl,anchor=west] at (9.2,2.15) {disputed: see AGI slide};
  \node[lbl,anchor=west] at (9.2,3.7) {hypothetical};
\end{tikzpicture}
\end{sfig}
{\small This ladder is a \textbf{communication device}, not a scientific classification. Modern LLMs blur it: one system handles many tasks but fails others a child finds easy.}
\discuss{Is ``general'' a yes/no property or a matter of degree?}
\end{frame}

% ---------- 2. History ----------
\begin{frame}{A short history of AI, 1943 to 2026}
\begin{sfig}[0.78]
\begin{tikzpicture}[x=1mm,y=1mm]
  \draw[very thick,tdMuted] (0,0) -- (152,0);
  \foreach \y in {1950,1960,1970,1980,1990,2000,2010,2020}{\pgfmathsetmacro{\xx}{\y<2010 ? (\y-1940)*1.1 : 77+(\y-2010)*4.4}\draw[tdMuted] (\xx,-1.2) -- (\xx,1.2); \node[font=\tiny,text=tdMuted] at (\xx,-3.5) {\y};}
  \fill[tdRed!18] ({(1974-1940)*1.1},-1.4) rectangle ({(1980-1940)*1.1},1.4);
  \fill[tdRed!18] ({(1987-1940)*1.1},-1.4) rectangle ({(1995-1940)*1.1},1.4);
  \node[font=\tiny,text=tdRed] at ({(1977-1940)*1.1},-7) {winter};
  \node[font=\tiny,text=tdRed] at ({(1991-1940)*1.1},-7) {winter};
  \foreach \y/\t/\c/\h in {1943/{McCulloch--Pitts\\ neuron}/tdBlue/10,1950/{Turing:\\ imitation game}/tdAmber/24,1956/{Dartmouth\\ workshop}/tdGreen/10,1958/{Perceptron}/tdBlue/-14,1966/{ELIZA}/tdAmber/14,1997/{Deep Blue\\ beats Kasparov}/tdGreen/14,2012/{AlexNet,\\ deep learning}/tdBlue/12,2016/{AlphaGo}/tdGreen/-12,2017/{Transformer}/tdPurple/28,2020/{GPT-3}/tdPurple/-26,2022/{ChatGPT}/tdRed/12,2024/{reasoning\\ models}/tdAmber/-12,2025/{agents}/tdBlue/-26}{
    \pgfmathsetmacro{\xx}{\y<2010 ? (\y-1940)*1.1 : 77+(\y-2010)*4.4}
    \fill[\c] (\xx,0) circle (1.1mm);
    \ifnum\h<0 \draw[\c!60] (\xx,0) -- (\xx,{\h+5}); \node[font=\scriptsize,align=center,text=\c!70!black,anchor=north] at (\xx,{\h+5}) {\textbf{\y}\\ \t};
    \else \draw[\c!60] (\xx,0) -- (\xx,{\h-1}); \node[font=\scriptsize,align=center,text=\c!70!black,anchor=south] at (\xx,{\h-1}) {\textbf{\y}\\ \t};\fi
  }
\end{tikzpicture}
\end{sfig}
\vspace{-1mm}
{\scriptsize Horizontal scale compressed before 2010 and stretched after. Sources: McCulloch and Pitts 1943; Turing 1950; Dartmouth 1956; Rosenblatt 1958; Weizenbaum 1966; Deep Blue 1997; AlexNet 2012; AlphaGo 2016; Vaswani et al.\ 2017; GPT-3 2020; ChatGPT 30 Nov 2022. Reasoning models (2024) and agents (2025): see the ``AI now'' part. Winter dates approximate.}
\end{frame}

\begin{frame}{Dartmouth, summer 1956: the field gets its name}
\begin{columns}[T]
\begin{column}{0.52\textwidth}
\small
John McCarthy, Marvin Minsky, Nathaniel Rochester and Claude Shannon proposed a two-month study at Dartmouth College. The proposal introduced the phrase \textbf{``artificial intelligence''}.

The conjecture at its centre:
\vspace{1mm}

\qcard{tdBlue}{\linewidth-10pt}{``every aspect of learning or any other feature of intelligence can in principle be so precisely described that a machine can be made to simulate it.''}
\vspace{1mm}

They expected a ``significant advance'' from a small group over one summer.
\end{column}
\begin{column}{0.45\textwidth}
\begin{sfig}[0.46]\begin{tikzpicture}
  \node[card=tdBlue,text width=3.0cm,minimum height=1.8cm] (a) at (0,1.2) {\textbf{Symbolic AI}\\{\footnotesize logic, search, expert systems}};
  \node[card=tdGreen,text width=3.0cm,minimum height=1.8cm] (b) at (0,-1.2) {\textbf{Connectionism}\\{\footnotesize neural networks that learn}};
  \node[card=tdRed,text width=2.6cm,minimum height=1.2cm] (c) at (4.4,1.2) {\footnotesize two winters};
  \node[card=tdPurple,text width=2.6cm,minimum height=1.2cm] (d) at (4.4,-1.2) {\footnotesize deep learning and LLMs};
  \draw[arr] (a) -- (c); \draw[arr] (b) -- (d);
\end{tikzpicture}\end{sfig}
{\small The two traditions took turns being fashionable, and the current systems are connectionist at heart.}
\end{column}
\end{columns}
\src{McCarthy, Minsky, Rochester, Shannon, ``A Proposal for the Dartmouth Summer Research Project on Artificial Intelligence'', 31 Aug 1955.}
\end{frame}

\begin{frame}{Hype, winters and comebacks}
\begin{columns}[T]
\begin{column}{0.56\textwidth}
\begin{sfig}[0.62]
\begin{tikzpicture}
  \draw[->] (0,0) -- (8.6,0) node[right,font=\scriptsize] {time};
  \draw[->] (0,0) -- (0,3.4) node[above,font=\scriptsize] {expectations};
  \draw[very thick,tdBlue,smooth] plot coordinates {(0.2,0.4) (0.9,2.2) (1.6,3.0) (2.4,1.2) (3.0,0.6) (3.6,1.8) (4.3,2.7) (5.1,0.7) (6.0,0.5) (6.8,1.5) (7.4,2.6) (8.2,3.2)};
  \node[font=\tiny,align=center,anchor=south] at (1.6,3.0) {1950s--60s\\ first boom};
  \node[font=\tiny,align=center,text=tdRed,anchor=north] at (3.0,0.55) {first\\ winter};
  \node[font=\tiny,align=center,anchor=south] at (4.3,2.7) {1980s expert\\ systems};
  \node[font=\tiny,align=center,text=tdRed,anchor=north] at (5.8,0.45) {second\\ winter};
  \node[font=\tiny,align=center,anchor=south] at (8.2,3.2) {deep learning\\ and LLMs};
\end{tikzpicture}
\end{sfig}
{\scriptsize\illus}
\end{column}
\begin{column}{0.41\textwidth}
\small
\textbf{Why winters happened}
\begin{itemize}\setlength\itemsep{2pt}
\item promises ran ahead of results;
\item early methods did not scale (combinatorial explosion, brittle rules);
\item funders cut budgets, notably after the 1973 Lighthill report in the UK.
\end{itemize}
\textbf{What ended them:} more data, more compute, better learning methods.
\end{column}
\end{columns}
\discuss{We are in a boom. Which of the winter warning signs would you watch for today?}
\src{Lighthill, ``Artificial Intelligence: A General Survey'' (1973); Russell and Norvig, AIMA, 4th ed., ch.\ 1 (history).}
\end{frame}

\begin{frame}{The deep-learning turn}
\begin{columns}[T]
\begin{column}{0.5\textwidth}
\small
\begin{itemize}\setlength\itemsep{4pt}
\item[\key{2012}] \textbf{AlexNet} wins the ImageNet image-recognition contest by a wide margin, using GPUs and a deep convolutional network.
\item[\key{2016}] \textbf{AlphaGo} beats Lee Sedol 4--1 at Go, a game once thought decades away.
\item[\key{2017}] The \textbf{Transformer} (``Attention Is All You Need'') replaces recurrence with attention and scales well.
\item[\key{2020}] \textbf{GPT-3}: 175 billion parameters, few-shot learning from a prompt.
\item[\key{2022}] \textbf{ChatGPT} (30 Nov): a chat interface brings LLMs to hundreds of millions of people.
\end{itemize}
\end{column}
\begin{column}{0.47\textwidth}
\begin{sfig}[0.42]
\begin{tikzpicture}
\begin{axis}[width=7.4cm,height=4.4cm,ymode=log,xmin=2011,xmax=2021,ymin=2e7,ymax=1e12,xlabel={year},ylabel={parameters},xtick={2012,2017,2019,2020},xticklabels={2012,2017,2019,2020},ymajorgrids,grid style={tdLine}]
\addplot[only marks,mark=*,mark size=2.4pt,tdBlue] coordinates {(2012,6e7) (2017,6.5e7) (2019,1.5e9) (2020,1.75e11)};
\node[font=\tiny,anchor=south west] at (axis cs:2012.1,8e7) {AlexNet 60M};
\node[font=\tiny,anchor=north west] at (axis cs:2017.1,5e7) {Transformer (base) 65M};
\node[font=\tiny,anchor=south east] at (axis cs:2018.9,2.2e9) {GPT-2 1.5B};
\node[font=\tiny,anchor=north east] at (axis cs:2019.9,1.3e11) {GPT-3 175B};
\end{axis}
\end{tikzpicture}
\end{sfig}
{\scriptsize Parameter counts as reported in the original papers; recipe: big Transformer, next-token prediction, human feedback, scale.}
{\small The lesson many researchers took: \textbf{general methods that scale} beat hand-written knowledge (Sutton's ``bitter lesson'', 2019).}
\end{column}
\end{columns}
\src{Krizhevsky et al., NeurIPS 2012; Silver et al., Nature 2016; Vaswani et al., NeurIPS 2017; Brown et al., NeurIPS 2020; OpenAI, ChatGPT launch 30 Nov 2022; Sutton, ``The Bitter Lesson'' (2019).}
\end{frame}

\begin{frame}{From chatbots to reasoning models to agents}
\begin{sfig}[0.45]
\begin{tikzpicture}
  \node[card=tdBlue,minimum width=3.4cm,minimum height=2cm] (a) at (0,0) {\textbf{Pretrained LLM}\\{\scriptsize predicts next token}};
  \node[card=tdGreen,minimum width=3.4cm,minimum height=2cm] (b) at (4.3,0) {\textbf{Chat assistant}\\\shortstack{\scriptsize instructions +\\ human feedback}};
  \node[card=tdAmber,minimum width=3.4cm,minimum height=2cm] (c) at (8.6,0) {\textbf{Reasoning model}\\\shortstack{\scriptsize ``thinks'' longer at\\ test time}};
  \node[card=tdRed,minimum width=3.4cm,minimum height=2cm] (d) at (12.9,0) {\textbf{Agent}\\\shortstack{\scriptsize tools, code, browser,\\ long tasks}};
  \draw[arr] (a)--(b); \draw[arr] (b)--(c); \draw[arr] (c)--(d);
  \node[lbl] at (0,-1.6) {2018--2021}; \node[lbl] at (4.3,-1.6) {2022}; \node[lbl] at (8.6,-1.6) {2024--2025}; \node[lbl] at (12.9,-1.6) {2025--2026};
\end{tikzpicture}
\end{sfig}
\vspace{1mm}
\small
The ARC Prize 2025 analysis calls \textbf{refinement loops} (generate a candidate, check it against feedback, try again) ``the central theme'' of 2025 progress: the same idea behind reasoning models and agents, spending compute at \emph{inference} rather than only at training.
\discuss{If a model improves by ``thinking longer'', is that more intelligence, or more compute?}
\src{ARC Prize Foundation, ``ARC Prize 2025 Results and Analysis'', arcprize.org/blog/arc-prize-2025-results-analysis (Dec 2025); date ranges for the stages are approximate.}
\end{frame}

% ---------- 3. Turing test ----------
\begin{frame}{Turing 1950: replace ``can machines think?'' by a game}
\begin{columns}[T]
\begin{column}{0.48\textwidth}
\small
In ``Computing Machinery and Intelligence'' (\emph{Mind}, 1950) Turing opens: ``I propose to consider the question, `Can machines think?'\,'' and then swaps it for a \textbf{behavioural test}, the \emph{imitation game}.

\begin{itemize}\setlength\itemsep{3pt}
\item Original game: interrogator C must tell man A from woman B by text; A tries to deceive.
\item Then a \textbf{machine} takes A's part.
\item Is C wrong as often as before?
\end{itemize}
\vspace{1mm}
\qcard{tdAmber}{\linewidth-10pt}{\scriptsize\textbf{Turing's prediction:} in about fifty years computers would play so well that an average interrogator would have no more than a \textbf{70\% chance} of identifying the machine after \textbf{five minutes} of questioning.}
\end{column}
\begin{column}{0.49\textwidth}
\begin{sfig}[0.62]
\begin{tikzpicture}
  \node[card=tdBlue,minimum width=1.7cm] (c) at (0,0) {\textbf{C}\\{\scriptsize interrogator}};
  \node[card=tdGreen,minimum width=1.7cm] (a) at (4.6,1.1) {\textbf{A}\\{\scriptsize man, deceives}};
  \node[card=tdGreen,minimum width=1.7cm] (b) at (4.6,-1.1) {\textbf{B}\\{\scriptsize woman, helps}};
  \draw[<->,thick,tdMuted] (c) -- node[above,lbl,sloped] {typed text} (a);
  \draw[<->,thick,tdMuted] (c) -- (b);
  \node[card=tdRed,minimum width=1.7cm] (m) at (8.2,1.1) {\textbf{machine}\\{\scriptsize takes A's role}};
  \draw[arr,tdRed] (m) -- node[above,lbl] {replaces} (a);
  \node[lbl] at (2.3,-2.3) {C sees only the text, never the players};
\end{tikzpicture}
\end{sfig}
\end{column}
\end{columns}
\src{A.\ M.\ Turing, ``Computing Machinery and Intelligence'', Mind 59(236), 433--460, 1950.}
\end{frame}

\begin{frame}{The standard three-party Turing test}
\begin{sfig}[0.55]
\begin{tikzpicture}
  \node[card=tdBlue,minimum width=2.4cm,minimum height=1.6cm] (j) at (0,0) {\textbf{Judge}\\\shortstack{\scriptsize chats with both\\ at the same time}};
  \node[card=tdGreen,minimum width=2.4cm,minimum height=1.2cm] (h) at (6,1.4) {\textbf{Human}\\{\scriptsize tries to seem human}};
  \node[card=tdRed,minimum width=2.4cm,minimum height=1.2cm] (m) at (6,-1.4) {\textbf{AI system}\\{\scriptsize tries to seem human}};
  \draw[<->,thick,tdMuted] (j) -- node[above,lbl,sloped] {5-minute text chat} (h);
  \draw[<->,thick,tdMuted] (j) -- (m);
  \node[pill=tdAmber] (v) at (11,0) {verdict: ``which one is the human?''};
  \draw[arr] (j) .. controls (3,-3.4) and (9,-3.4) .. (v.south);
  \node[lbl,anchor=west] at (7.5,-2.9) {AI ``wins'' if judged human more than 50\% of the time};
\end{tikzpicture}
\end{sfig}
\vspace{1mm}
{\small Three-party tests are \textbf{stricter} than two-party ones: the judge can compare the two partners directly, so the guess is harder to get right by luck. Chance level is \textbf{50\%}.}
\discuss{Why might the judge find a human \emph{less} human than a machine? What would you ask in five minutes?}
\end{frame}

\begin{frame}{ELIZA (1966): a mirror that people mistook for a mind}
\begin{columns}[T]
\begin{column}{0.5\textwidth}
\small
Joseph Weizenbaum (MIT) built \textbf{ELIZA}; its DOCTOR script imitated a Rogerian therapist who reflects the patient's words back as questions.
\begin{itemize}\setlength\itemsep{3pt}
\item no model of the world;
\item keywords + templates + pronoun swapping;
\item yet some users confided in it.
\end{itemize}
Weizenbaum was disturbed by how readily people credited it with understanding: the \textbf{ELIZA effect}.
\end{column}
\begin{column}{0.47\textwidth}
\begin{sfig}[0.55]
\begin{tikzpicture}
  \node[soft,text width=5.3cm,align=left] (u) at (0,1.5) {\textbf{User:} I am sad about my exam};
  \node[card=tdAmber,text width=5.3cm,align=left] (r) at (0,0) {\scriptsize rule \texttt{i am (.*)}\\ swap pronouns: my $\to$ your\\ fill template: \texttt{Why are you \{\}?}};
  \node[soft,text width=5.3cm,align=left] (b) at (0,-1.6) {\textbf{ELIZA:} Why are you sad about your exam?};
  \draw[arr] (u) -- (r); \draw[arr] (r) -- (b);
\end{tikzpicture}
\end{sfig}
\end{column}
\end{columns}
\discuss{People are eager to see minds in text. Is that a flaw in us, or a feature that machines can exploit?}
\src{Weizenbaum, ``ELIZA: a computer program for the study of natural language communication between man and machine'', CACM 9(1), 1966; Weizenbaum, \emph{Computer Power and Human Reason}, 1976.}
\end{frame}

\begin{frame}[fragile]{Demo 1: ELIZA in 20 lines of Python 3.14}
\begin{columns}[T]
\begin{column}{0.58\textwidth}
\pycodes{eliza_eliza}
\end{column}
\begin{column}{0.4\textwidth}
\small
\textbf{How it works}
\begin{itemize}\setlength\itemsep{2pt}
\item rules are tried top to bottom; the first regex that \emph{fully matches} wins;
\item \texttt{reflect} turns ``my'' into ``your'', ``me'' into ``you'';
\item the last rule matches anything, so ELIZA \textbf{never fails}, it just stalls.
\end{itemize}
{\footnotesize Nothing here represents \emph{sadness}, \emph{exams} or \emph{families}: they are just strings.}
\end{column}
\end{columns}
\end{frame}

\begin{frame}[fragile]{Demo 1 output: it feels attentive, until it does not}
\begin{columns}[T]
\begin{column}{0.56\textwidth}
\pyouts{eliza_talk}
\end{column}
\begin{column}{0.42\textwidth}
\small
\begin{itemize}\setlength\itemsep{3pt}
\item The first four replies sound \textbf{engaged}: pure string surgery.
\item The last one shows the seam: asked about itself, it falls through to the catch-all.
\end{itemize}
\vspace{1mm}
\qcard{tdBlue}{\linewidth-10pt}{\footnotesize \textbf{Lesson.} Passing a short conversation can be done with shallow tricks. That is why the \emph{design} of a test (length, judge training, adversarial questions) matters as much as the result.}
\end{column}
\end{columns}
\end{frame}

\begin{frame}{Loebner Prize and ``Eugene Goostman''}
\begin{columns}[T]
\begin{column}{0.5\textwidth}
\small
\textbf{Loebner Prize.} An annual chatbot contest run from 1991 on a Turing-style format. Winners typically relied on \emph{tricks}: deflecting, changing the subject, typos, humour.

\vspace{2mm}
\textbf{Eugene Goostman (2014).} At a Royal Society event organised by the University of Reading, a chatbot posing as a 13-year-old Ukrainian boy was reported to have fooled 33\% of judges (threshold 30\%) and was widely announced as passing the test.

\vspace{2mm}
\begin{sfig}[0.2]\begin{tikzpicture}[x=0.07cm]
  \draw[tdMuted] (0,0) -- (100,0);
  \fill[tdBlue] (0,0.15) rectangle (33,0.65);
  \draw[tdRed,thick] (30,-0.1) -- (30,0.9);
  \node[font=\scriptsize,anchor=west] at (35,0.4) {33\% of judges fooled};
  \node[font=\scriptsize,text=tdRed,anchor=south] at (30,0.9) {30\% threshold};
\end{tikzpicture}\end{sfig}
\end{column}
\begin{column}{0.47\textwidth}
\begin{block}{Why many experts were not convinced}
{\small
\begin{itemize}\setlength\itemsep{2pt}
\item a persona (young, non-native English) \textbf{excuses} errors and evasions;
\item only five-minute conversations;
\item the 30\% figure is a reading of Turing's remark, not a rule he gave;
\item the bot used the same tricks as ELIZA with a bigger script.
\end{itemize}}
\end{block}
\end{column}
\end{columns}
\vspace{1mm}
{\small \textbf{Persona and context are levers.} Keep this in mind for the 2025 results.}
\src{Loebner Prize history: loebner.net / AISB; University of Reading press release, 8 June 2014 (``Turing test success marks milestone in computing history'').}
\end{frame}

\begin{frame}{Updated results: three-party Turing test (Jones and Bergen)}
\begin{columns}[T]
\begin{column}{0.6\textwidth}
\begin{sfig}[0.75]
\begin{tikzpicture}
\begin{axis}[xbar,width=9.4cm,height=6.0cm,xmin=0,xmax=100,xlabel={judged human by the interrogator (\%)},symbolic y coords={GPT-4o,ELIZA,GPT-4.5 plain,LLaMa-3.1 plain,LLaMa-3.1 persona,GPT-4.5 persona},ytick=data,bar width=7pt,nodes near coords,nodes near coords align={horizontal},every node near coord/.append style={font=\scriptsize},tick label style={font=\scriptsize},xmajorgrids,grid style={tdLine},enlarge y limits=0.12]
\addplot[fill=tdMuted!50,draw=none] coordinates {(21,GPT-4o) (23,ELIZA) (36,GPT-4.5 plain) (38,LLaMa-3.1 plain)};
\addplot[fill=tdBlue,draw=none] coordinates {(56,LLaMa-3.1 persona) (73,GPT-4.5 persona)};
\draw[tdRed,dashed,thick] (rel axis cs:0.5,0) -- (rel axis cs:0.5,1);
\node[font=\tiny,text=tdRed,anchor=south west] at (rel axis cs:0.51,0.01) {chance = 50\%};
\end{axis}
\end{tikzpicture}
\end{sfig}
\end{column}
\begin{column}{0.37\textwidth}
\small
\textbf{Design:} 5-minute chats, one interrogator, one human and one AI at once; two independent groups (UC San Diego undergraduates, Prolific workers), nearly 500 participants.

\textbf{Result:} with a detailed human \emph{persona prompt}, GPT-4.5 was picked as the human \textbf{73\%} of the time, more often than the real human. LLaMa-3.1-405B: 56\%. Without the persona: 36\% and 38\%.

\textbf{Venue:} PNAS, 2026 (preprint March 2025).
\end{column}
\end{columns}
\src{Jones and Bergen, ``Large language models pass a standard three-party Turing test'', PNAS 2026, doi:10.1073/pnas.2524472123; preprint arXiv:2503.23674 (31 Mar 2025); UC San Diego news release, 19 May 2026 (reported via techxplore.com/news/2026-05-ai-human-real-humans-classic.html). Figures are those reported in the release.}
\end{frame}

\begin{frame}{How to read the result (and what came before)}
\begin{columns}[T]
\begin{column}{0.5\textwidth}
\small
\textbf{What it shows}
\begin{itemize}\setlength\itemsep{3pt}
\item In short, text-only, adversarial-free chats a prompted LLM can be \textbf{indistinguishable from, or more ``human'' than}, a human.
\item Earlier (2024, two-party): GPT-4 \textbf{54\%}, ELIZA 22\%, humans 67\%.
\end{itemize}
\textbf{What it does not show}
\begin{itemize}\setlength\itemsep{3pt}
\item Without the persona: only 36--38\%.
\item 5 (or 15) minutes, non-expert judges.
\item Not evidence of understanding or consciousness.
\end{itemize}
\end{column}
\begin{column}{0.47\textwidth}
\qcard{tdAmber}{\linewidth-10pt}{\small\textbf{The authors' own reading}\\[2pt]
Jones: ``The Turing test is a game about lying for the models.''\\[2pt]
Bergen: ``Increasingly, it's measuring humanlikeness.''\\[2pt]
\scriptsize UC San Diego news release, May 2026.}
\vspace{2mm}
\par\small The models won partly by \emph{human-like mistakes}: informal style, typos, not knowing things.
\end{column}
\end{columns}
\discuss{If a machine must dumb itself down to pass, is that intelligence or acting?}
\src{Jones and Bergen, ``People cannot distinguish GPT-4 from a human in a Turing test'', arXiv:2405.08007 (2024); UC San Diego news release 19 May 2026 (via techxplore.com).}
\end{frame}

\begin{frame}{Critiques of the Turing test}
\begin{sfig}[0.62]
\begin{tikzpicture}
  \node[card=tdBlue,text width=4.6cm,minimum height=2.4cm] at (0,0) {\textbf{Tests humanlikeness, not intelligence}\\{\scriptsize A smart system may be easy to spot; a shallow one may fool people (ELIZA, Goostman).}};
  \node[card=tdGreen,text width=4.6cm,minimum height=2.4cm] at (5.2,0) {\textbf{Rewards deception}\\{\scriptsize The machine must lie about being a machine, and about its nature, tastes and body.}};
  \node[card=tdAmber,text width=4.6cm,minimum height=2.4cm] at (10.4,0) {\textbf{Anthropocentric}\\{\scriptsize ``Submarines swim'' objection: attributed to Dijkstra, the question is no more useful than whether submarines swim.}};
  \node[card=tdRed,text width=4.6cm,minimum height=2.4cm] at (0,-3) {\textbf{Behaviour is not the whole story}\\{\scriptsize Searle's Chinese Room: perfect output without understanding (next part).}};
  \node[card=tdPurple,text width=4.6cm,minimum height=2.4cm] at (5.2,-3) {\textbf{Depends on judge and setup}\\{\scriptsize Duration, expertise, allowed topics and incentives change the outcome.}};
  \node[card=tdBlue,text width=4.6cm,minimum height=2.4cm] at (10.4,-3) {\textbf{Defence}\\{\scriptsize Turing meant it as a \emph{sufficient} behavioural criterion and a way out of vague definitions, not as the definition of intelligence.}};
\end{tikzpicture}
\end{sfig}
\discuss{Design a better test of intelligence in one sentence. What loophole would a clever machine find?}
\src{Turing (1950); Searle, ``Minds, Brains, and Programs'', Behavioral and Brain Sciences 3(3), 1980; Hayes and Ford, ``Turing Test Considered Harmful'', IJCAI 1995; Russell and Norvig, AIMA 4th ed., ch.\ 1.}
\end{frame}

% ---------- 4. Searle and understanding ----------
\begin{frame}{Searle's Chinese Room (1980)}
\begin{columns}[T]
\begin{column}{0.36\textwidth}
\small
\textbf{Thought experiment.} You, who understand no Chinese, sit in a closed room with a big rule book (in English) and baskets of Chinese symbols.

Slips with Chinese questions come in. You follow the rules, match shapes, and pass out other slips.

To the people outside, your answers are \textbf{fluent Chinese}. Inside, you understand \textbf{nothing}.

\vspace{1mm}
\textbf{Searle's point:} a program only manipulates symbols by their \emph{form} (syntax). That cannot, by itself, give \emph{meaning} (semantics).
\end{column}
\begin{column}{0.61\textwidth}
\begin{sfig}[0.7]
\begin{tikzpicture}[scale=1.0,every node/.append style={font=\small}]
  \draw[rounded corners=6pt,thick,draw=tdInk!70,fill=tdPaper!60!tdLine!30] (2.2,-1.9) rectangle (8.2,1.9);
  \node[font=\scriptsize\bfseries,anchor=north west] at (2.3,1.85) {the room};
  \node[card=tdBlue,minimum width=1.6cm] (p) at (4.0,0) {\textbf{you}\\{\scriptsize no Chinese}};
  \node[card=tdAmber,minimum width=1.7cm] (r) at (6.6,0.7) {\textbf{rule book}\\\shortstack{\scriptsize if shape X,\\ send shape Y}};
  \node[card=tdGreen,minimum width=1.7cm] (s) at (6.6,-0.9) {\textbf{symbols}\\{\scriptsize baskets}};
  \draw[arr] (p) -- (r); \draw[arr] (p) -- (s);
  \node[soft,minimum width=1.5cm] (in) at (-0.4,0.8) {\scriptsize\ttfamily ni hao ma?};
  \node[soft,minimum width=1.5cm] (out) at (10.5,-0.8) {\scriptsize\ttfamily wo hen hao};
  \draw[arr] (in) -- (2.2,0.4); \draw[arr] (8.2,-0.4) -- (out);
  \node[lbl] at (-0.4,1.6) {question in}; \node[lbl] at (10.5,0.0) {answer out};
\end{tikzpicture}
\end{sfig}
\end{column}
\end{columns}
\src{Searle, ``Minds, Brains, and Programs'', Behavioral and Brain Sciences 3(3), 417--457, 1980.}
\end{frame}

\begin{frame}{The argument, step by step}
\begin{columns}[T]
\begin{column}{0.55\textwidth}
\small
\begin{enumerate}\setlength\itemsep{4pt}
\item A computer program is purely \textbf{syntactic}: it manipulates symbols by their form.
\item Minds have \textbf{semantics}: thoughts are \emph{about} things.
\item Syntax is not sufficient for semantics (the room shows this).
\item \textbf{So} running a program is not sufficient for having a mind.
\end{enumerate}
\vspace{1mm}
Searle's targets: \textbf{strong AI} (the right program \emph{is} a mind) versus \textbf{weak AI} (programs are useful tools for studying or simulating minds). He accepts weak AI.

{\footnotesize He does \emph{not} say machines cannot think: brains are physical machines that cause minds. He says the \emph{causal powers of brains} matter, not just the program.}
\end{column}
\begin{column}{0.42\textwidth}
\begin{sfig}[0.6]
\begin{tikzpicture}
  \node[card=tdBlue,minimum width=4cm] (a) at (0,2.2) {\textbf{Programs}\\{\scriptsize syntax only}};
  \node[card=tdGreen,minimum width=4cm] (b) at (0,0.2) {\textbf{Minds}\\{\scriptsize have meaning (semantics)}};
  \node[card=tdRed,minimum width=4cm] (c) at (0,-2.0) {\textbf{Conclusion}\\\shortstack{\scriptsize programs alone do\\ not suffice for minds}};
  \draw[arr] (a) -- node[right,lbl] {syntax $\neq$ semantics} (c);
  \draw[arr] (b) -- (c);
\end{tikzpicture}
\end{sfig}
\end{column}
\end{columns}
\end{frame}

\begin{frame}[fragile]{Demo 2: a Chinese Room in a dozen lines}
\begin{columns}[T]
\begin{column}{0.5\textwidth}
\pycodes{chinese_room_room}
\vspace{1mm}
{\footnotesize The rule book is a Python dictionary: 4 rules, written in romanised Chinese.}
\end{column}
\begin{column}{0.47\textwidth}
\pyouts{chinese_room_run}
\vspace{1mm}
\pyouts{chinese_room_stats}
\vspace{1mm}
{\footnotesize On the four known slips the room \emph{looks} competent. On a new question it stalls with ``please say that again''. A real room would need an astronomically large book, which is exactly what the replies below argue about.}
\end{column}
\end{columns}
\end{frame}

\begin{frame}{Four classic replies, and Searle's answers}
\begin{sfig}[0.66]
\begin{tikzpicture}
  \node[card=tdBlue,text width=6.4cm,minimum height=2.8cm] at (0,0) {\textbf{Systems reply}\\{\footnotesize The \emph{whole system} (you + book + room) understands, not the man alone.}\\[3pt]{\scriptsize\color{tdRed!70!black}Searle: let the man memorise everything; still no understanding.}};
  \node[card=tdGreen,text width=6.4cm,minimum height=2.8cm] at (7.2,0) {\textbf{Robot reply}\\{\footnotesize Put the program in a robot with cameras and limbs: symbols get \emph{grounded} in the world.}\\[3pt]{\scriptsize\color{tdRed!70!black}Searle: camera input is more symbols; the man still just shuffles them.}};
  \node[card=tdAmber,text width=6.4cm,minimum height=2.8cm] at (0,-3.3) {\textbf{Brain simulator reply}\\{\footnotesize Simulate every neuron of a Chinese speaker's brain.}\\[3pt]{\scriptsize\color{tdRed!70!black}Searle: simulating a storm does not make you wet; simulation is not duplication.}};
  \node[card=tdPurple,text width=6.4cm,minimum height=2.8cm] at (7.2,-3.3) {\textbf{Other minds reply}\\{\footnotesize We attribute understanding to people only from behaviour. Same standard for machines?}\\[3pt]{\scriptsize\color{tdRed!70!black}Searle: the issue is what understanding \emph{is}, not how we find out about it.}};
\end{tikzpicture}
\end{sfig}
\discuss{Which reply do you find strongest? Which of Searle's answers feels like a dodge?}
\src{Searle (1980) and the open peer commentary printed with it; Stanford Encyclopedia of Philosophy, ``The Chinese Room Argument''.}
\end{frame}

\begin{frame}{Do large language models understand?}
\begin{sfig}[0.55]
\begin{tikzpicture}
  \node[card=tdRed,text width=7.0cm,minimum height=3.7cm,anchor=north east] at (-0.2,0) {\textbf{``Stochastic parrots''}\\[2pt]\footnotesize Bender et al.\ (2021): LMs stitch together text by probability, without reference to meaning. Bender and Koller (2020): form alone cannot teach meaning (the ``octopus test'').\\[2pt]{\scriptsize\color{tdMuted} Warns of bias, cost, and people mistaking fluency for a mind.}};
  \node[card=tdGreen,text width=7.0cm,minimum height=3.7cm,anchor=north west] at (0.2,0) {\textbf{``Something more is going on''}\\[2pt]\footnotesize LLMs solve novel problems, use abstractions and show internal structure: world-model-like features (board-game and spatial probes), planning in circuits. Prediction at scale may force useful internal models.\\[2pt]{\scriptsize\color{tdMuted} Evidence is partial and contested.}};
  \node[pill=tdAmber] at (0,-4.2) {middle positions exist: some competence, little grounding, unclear what ``understanding'' means};
\end{tikzpicture}
\end{sfig}
\discuss{What observable evidence would convince you that a model understands? What would convince you it does not?}
\src{Bender and Koller, ACL 2020; Bender, Gebru, McMillan-Major, Shmitchell, ``On the Dangers of Stochastic Parrots'', FAccT 2021; Li et al., ``Emergent World Representations'', ICLR 2023.}
\end{frame}

\begin{frame}[fragile]{Demo 3: a tiny ``parrot''}
\begin{columns}[T]
\begin{column}{0.52\textwidth}
\pycodes{parrot_parrot}
\end{column}
\begin{column}{0.45\textwidth}
\pyouts{parrot_run}
\vspace{1mm}
{\footnotesize
Each word is chosen using \textbf{only the previous word}. The text is locally plausible and globally empty.

\textbf{Careful:} an LLM is \emph{not} a bigram table. It conditions on thousands of earlier tokens and has billions of learned parameters. The demo shows what ``predicting the next word'' alone does \emph{not} prove: that fluency needs no understanding at tiny scale. Whether that stays true at scale is the debate.}
\end{column}
\end{columns}
\end{frame}

\begin{frame}{Emergent abilities: real or a mirage?}
\begin{columns}[T]
\begin{column}{0.5\textwidth}
\small
\textbf{Claim (Wei et al.\ 2022).} Some abilities (multi-step arithmetic, certain reasoning tasks) are \textbf{absent in small models and appear abruptly} past a scale threshold.

\textbf{Challenge (Schaeffer et al.\ 2023).} Much of the ``jump'' comes from the \textbf{metric}: all-or-nothing scores (exact match) hide smooth underlying progress. Use a continuous metric and the cliff often becomes a slope.

\vspace{1mm}
Take-away: be careful with stories about sudden ``sparks''. Ask what was measured.
\end{column}
\begin{column}{0.47\textwidth}
\begin{sfig}[0.6]
\begin{tikzpicture}
\begin{axis}[width=7.4cm,height=5.2cm,xlabel={model scale (log)},ylabel={task score},ymin=-0.05,ymax=1.1,xmin=0,xmax=10,xtick=\empty,ytick={0,0.5,1},legend pos=north west]
\addplot[very thick,tdRed,domain=0:10,samples=60] {1/(1+exp(-(x-6.5)*3.0))};
\addplot[very thick,tdGreen,domain=0:10,samples=60] {0.05+0.9*x/10};
\legend{exact-match metric,continuous metric}
\end{axis}
\end{tikzpicture}
\end{sfig}
{\scriptsize\illus: the same hypothetical improvement, scored two ways.}
\end{column}
\end{columns}
\src{Wei et al., ``Emergent Abilities of Large Language Models'', TMLR 2022; Schaeffer, Miranda, Koyejo, ``Are Emergent Abilities of Large Language Models a Mirage?'', NeurIPS 2023.}
\end{frame}

\begin{frame}{Looking inside: mechanistic interpretability}
\begin{columns}[T]
\begin{column}{0.52\textwidth}
\small
Instead of arguing from outside, \textbf{open the box}: find the features and circuits a network uses.

Anthropic's March 2025 \emph{circuit tracing} study of Claude 3.5 Haiku reported, for example:
\begin{itemize}\setlength\itemsep{3pt}
\item when writing a rhyming couplet, the model \textbf{chooses the rhyme word before writing the line}; editing that internal ``plan'' changes the ending;
\item shared, language-independent concepts across languages;
\item in one hard maths case, a plausible-sounding explanation that \textbf{did not match} the computation actually done.
\end{itemize}
\end{column}
\begin{column}{0.45\textwidth}
\begin{sfig}[0.3]\begin{tikzpicture}
  \node[sk=tdBlue,text width=2.0cm] (a) at (0,0) {line 1 of a couplet ends};
  \node[sk=tdAmber,text width=2.2cm] (b) at (3.4,0) {internal feature: rhyme word chosen};
  \node[sk=tdGreen,text width=2.0cm] (c) at (6.8,0) {line 2 written toward it};
  \draw[arr] (a) -- (b); \draw[arr] (b) -- (c);
  \node[font=\tiny,text=tdRed!70!black] at (3.4,-1.0) {suppress ``rabbit'': ending changes};
\end{tikzpicture}\end{sfig}
\begin{block}{Honest limits (from the authors)}
{\small
\begin{itemize}\setlength\itemsep{2pt}
\item only a \textbf{fraction} of the computation;
\item hours of effort per \textbf{short prompt};
\item possible artefacts of the method.
\end{itemize}}
\end{block}
{\footnotesize Planning and internal representations count \emph{against} the ``pure parrot'' picture, but say nothing about \textbf{consciousness}.}
\end{column}
\end{columns}
\src{Anthropic, ``Tracing the thoughts of a large language model'' (27 Mar 2025) and ``On the Biology of a Large Language Model'', transformer-circuits.pub (2025); open-source tooling announced 29 May 2025.}
\end{frame}

% ---------- 5. Other tests ----------
\begin{frame}{Winograd schemas: common sense in one word}
\begin{columns}[T]
\begin{column}{0.55\textwidth}
\small
Levesque's alternative to the Turing test: \textbf{sentence pairs differing by one word} that flips who a pronoun refers to. Answering needs everyday knowledge, not keyword tricks.

\vspace{1mm}
\qcard{tdBlue}{\linewidth-10pt}{The trophy doesn't fit in the suitcase because \textbf{it} is too \textbf{big}. $\rightarrow$ \emph{trophy}\\[3pt]
The trophy doesn't fit in the suitcase because \textbf{it} is too \textbf{small}. $\rightarrow$ \emph{suitcase}}
\vspace{1mm}

\textbf{Status.} Fell to large models: GPT-3 (2020) reported roughly 88--90\% on the 273-problem challenge; WinoGrande (44k problems) was built to be harder. Resisting tricks is not resisting \emph{scale}.
\end{column}
\begin{column}{0.4\textwidth}
\begin{sfig}[0.55]
\begin{tikzpicture}
  \node[card=tdAmber,minimum width=2.5cm] (t) at (0,1) {trophy};
  \node[card=tdGreen,minimum width=2.5cm] (s) at (0,-1) {suitcase};
  \node[pill=tdRed] (p) at (4,0) {``it''};
  \draw[arr,tdAmber] (p) -- node[above,lbl,sloped] {too big} (t);
  \draw[arr,tdGreen] (p) -- node[below,lbl,sloped] {too small} (s);
\end{tikzpicture}
\end{sfig}
\end{column}
\end{columns}
\src{Levesque, Davis, Morgenstern, ``The Winograd Schema Challenge'', KR 2012; Brown et al., ``Language Models are Few-Shot Learners'', NeurIPS 2020; Sakaguchi et al., ``WinoGrande'', AAAI 2020.}
\end{frame}

\begin{frame}{Benchmarks: a life cycle}
\begin{sfig}[0.42]
\begin{tikzpicture}
  \foreach \i/\t/\c in {0/{New benchmark:\\ models score low}/tdBlue,1/{Labs optimise\\ for it}/tdGreen,2/{Scores climb,\\ maybe leak or overfit}/tdAmber,3/{Saturated:\\ no longer informative}/tdRed,4/{Harder\\ benchmark}/tdPurple}{
    \node[card=\c,minimum width=2.8cm,minimum height=1.6cm] (b\i) at ({\i*3.4},0) {\t};}
  \foreach \a/\b in {0/1,1/2,2/3,3/4}{\draw[arr] (b\a)--(b\b);}
  \draw[arr,dashed] (b4.south) -- ++(0,-0.9) -| (b0.south);
\end{tikzpicture}
\end{sfig}
\small
\textbf{Goodhart's law:} when a measure becomes a target, it stops being a good measure. The ARC Prize team itself notes it suspects public training data may be \emph{inflating} ARC-AGI-1 and ARC-AGI-2 scores.
\discuss{Is a saturated benchmark evidence of progress, of overfitting, or both? How would you tell?}
\src{ARC Prize Foundation, ``ARC Prize 2025 Results and Analysis'', arcprize.org/blog/arc-prize-2025-results-analysis (Dec 2025).}
\end{frame}

\begin{frame}{ARC-AGI: puzzles that cannot be memorised}
\begin{columns}[T]
\begin{column}{0.5\textwidth}
\small
Chollet's \textbf{Abstraction and Reasoning Corpus} (2019): small coloured grids. You see 2--3 input/output examples, infer the \emph{rule}, and apply it to a new input. No language, no trivia, only priors a child has (objects, counting, symmetry).

\begin{itemize}\setlength\itemsep{2pt}
\item \textbf{ARC-AGI-1} (2019), \textbf{ARC-AGI-2} (Mar 2025): harder, resists brute force;
\item \textbf{ARC-AGI-3} (25 Mar 2026): interactive game-like environments (next slides).
\end{itemize}
\end{column}
\begin{column}{0.47\textwidth}
\begin{sfig}[0.6]
\begin{tikzpicture}[x=5.5mm,y=5.5mm]
  \tikzset{g/.style={draw=tdInk!50,minimum size=5.5mm,inner sep=0pt}}
  \foreach \c/\r in {0/0,1/1,2/2,3/3}{\fill[tdBlue] (\c,-\r) rectangle ++(1,-1);}
  \foreach \r in {0,...,3}{\foreach \c in {0,...,3}{\draw[tdInk!50] (\c,-\r) rectangle ++(1,-1);}}
  \node[font=\scriptsize,anchor=north] at (2,-4.2) {example input};
  \draw[arr] (4.5,-2) -- (6.3,-2);
  \begin{scope}[xshift=7cm]
  \foreach \c/\r in {0/0,1/1,2/2,3/3}{\fill[tdBlue] (\c,-\r) rectangle ++(1,-1);}
  \foreach \c/\r in {1/0,2/0,3/0,2/1,3/1,3/2}{\fill[tdGreen!60] (\c,-\r) rectangle ++(1,-1);}
  \foreach \r in {0,...,3}{\foreach \c in {0,...,3}{\draw[tdInk!50] (\c,-\r) rectangle ++(1,-1);}}
  \node[font=\scriptsize,anchor=north] at (2,-4.2) {example output};
  \end{scope}
\end{tikzpicture}
\end{sfig}
{\scriptsize Toy example in the ARC style (not an actual task): the rule is ``fill above the diagonal in green''. Real tasks are harder to put into words.}
\end{column}
\end{columns}
\src{Chollet, arXiv:1911.01547 (2019); Chollet et al., ``ARC-AGI-2'', arXiv:2505.11831 (2025); ARC Prize, ``Announcing ARC-AGI-3'', arcprize.org (25 Mar 2026).}
\end{frame}

\begin{frame}{ARC-AGI status, late November 2026}
\begin{columns}[T]
\begin{column}{0.55\textwidth}
\begin{sfig}[0.7]
\begin{tikzpicture}
\begin{axis}[ybar,width=9.2cm,height=6.0cm,ymin=0,ymax=118,ylabel={best score on leaderboard (\%)},symbolic x coords={ARC-AGI-1,ARC-AGI-2,ARC-AGI-3},xtick=data,bar width=18pt,nodes near coords,every node near coord/.append style={font=\scriptsize},ymajorgrids,grid style={tdLine},enlarge x limits=0.25]
\addplot[fill=tdBlue,draw=none] coordinates {(ARC-AGI-1,98.5) (ARC-AGI-2,95.8) (ARC-AGI-3,81.2)};
\draw[tdGreen,dashed,thick] (axis cs:ARC-AGI-1,100) -- (axis cs:ARC-AGI-3,100) node[right,font=\tiny,text=tdGreen!60!black] {};
\node[font=\tiny,text=tdGreen!60!black,anchor=south west] at (axis cs:ARC-AGI-1,100.5) {humans, ARC-AGI-3: 100\%};
\end{axis}
\end{tikzpicture}
\end{sfig}
\end{column}
\begin{column}{0.42\textwidth}
\small
Top entries listed on the ARC Prize leaderboard:
\begin{itemize}\setlength\itemsep{2pt}
\item ARC-AGI-1: 98.5\% (several models)
\item ARC-AGI-2: 95.8\%, GPT-6.2 Astra (17 Nov)
\item ARC-AGI-3: 81.2\%, Claude Fable 5.2 (9 Nov); GPT-6.2 Astra 70.3\%
\end{itemize}
\textbf{Caveats.} Runs differ in compute and harness (the top ARC-AGI-3 runs list roughly \$22--24k); ARC-AGI-1/2 are partly public. ARC Prize 2026 (Kaggle) has closed, results due 4 Dec 2026.
\end{column}
\end{columns}
\src{ARC Prize leaderboard, arcprize.org/leaderboard (read 23 Nov 2026; ARC-AGI-3 column is the ``Standard'' score); ARC Prize blog, arcprize.org/blog. Model names and dates are as listed there.}
\end{frame}

\begin{frame}{ARC-AGI-3 in 2026: from near zero to most of the way}
\begin{columns}[T]
\begin{column}{0.6\textwidth}
\begin{sfig}[0.72]
\begin{tikzpicture}
\begin{axis}[width=10.2cm,height=6.0cm,xmin=2,xmax=11,ymin=0,ymax=100,xlabel={2026},ylabel={best listed ARC-AGI-3 score (\%)},xtick={3,5,7,9,11},xticklabels={Mar,May,Jul,Sep,Nov},ymajorgrids,grid style={tdLine},ytick={0,25,50,75,100}]
\addplot[const plot,very thick,tdBlue,mark=*,mark size=1.8pt] coordinates {(2.8,0.5) (6.3,7.8) (6.8,30.2) (8.0,62.7) (10.3,81.2) (10.9,81.2)};
\node[font=\tiny,anchor=north west,align=left] at (axis cs:2.9,6) {launch, 25 Mar:\\ frontier $\approx$0.5\%};
\node[font=\tiny,anchor=south east,align=right] at (axis cs:6.2,10) {GPT-5.6 Sol\\ 7.8\% (9 Jul)};
\node[font=\tiny,anchor=east,align=right] at (axis cs:6.6,36) {Claude Opus 5\\ 30.2\% (24 Jul)};
\node[font=\tiny,anchor=east,align=right] at (axis cs:7.8,66) {GPT-6 Astra\\ 62.7\% (2 Sep)};
\node[font=\tiny,anchor=north,align=center] at (axis cs:9.6,76) {Claude Fable 5.2\\ 81.2\% (9 Nov)};
\end{axis}
\end{tikzpicture}
\end{sfig}
\end{column}
\begin{column}{0.37\textwidth}
\small
\textbf{What it tests.} Hundreds of hand-made, turn-based environments with \textbf{no instructions and no stated goal}: explore, infer the rules, work out how to win, carry knowledge to harder levels. Humans: 100\%.

\textbf{The surprise.} The launch post said frontier AI scored 0.51\%. Within eight months, the best listed entries moved to 70--81\%, at a high compute cost.

\textbf{Open question.} Does this mean the skill is solved, or that the benchmark is now being trained for?
\end{column}
\end{columns}
\src{arcprize.org/blog/arc-agi-3-launch (25 Mar 2026); leaderboard arcprize.org/leaderboard (read 23 Nov 2026). Points are the best listed entry at each date; scores come from different harnesses and budgets, so treat the curve as indicative.}
\end{frame}

\begin{frame}{Humanity's Last Exam (HLE)}
\begin{columns}[T]
\begin{column}{0.5\textwidth}
\small
\textbf{What.} 2,500 closed-ended, expert-written questions across more than a hundred subjects (multi-modal), built by the Center for AI Safety and Scale AI after older exams (MMLU) saturated. Published in \emph{Nature} on 28 Jan 2026.

\textbf{Why it matters.} It targets \textbf{knowledge and reasoning at the frontier of human expertise}, the opposite end from ARC-AGI's ``simple for humans'' puzzles.

\textbf{Scores.} Gemini 3 Deep Think: 48.4\% (press report, Feb 2026). Scale AI board, 17 Oct 2026: GPT-6 Astra 54.8\,$\pm$\,1.9\%. Aggregators list higher, partly self-reported, figures.
\end{column}
\begin{column}{0.47\textwidth}
\begin{sfig}[0.26]\begin{tikzpicture}[x=0.06cm]
  \draw[tdMuted] (0,0) -- (100,0);
  \fill[tdBlue!70] (0,1.0) rectangle (48.4,1.6); \node[font=\scriptsize,anchor=west] at (50,1.3) {Gemini 3 Deep Think 48.4\% (Feb)};
  \fill[tdBlue] (0,0.2) rectangle (54.8,0.8); \node[font=\scriptsize,anchor=west] at (56,0.5) {GPT-6 Astra 54.8\% (17 Oct)};
  \node[font=\tiny,text=tdMuted,anchor=north] at (50,-0.1) {score (\%), not like-for-like};
\end{tikzpicture}\end{sfig}
\begin{block}{Two kinds of ``hard''}
{\small
\begin{itemize}\setlength\itemsep{3pt}
\item \textbf{HLE:} hard for \emph{experts}.
\item \textbf{ARC-AGI:} easy for \emph{people}.
\item Leading on one says little about the other.
\end{itemize}}
\end{block}
\end{column}
\end{columns}
\src{Center for AI Safety, Scale AI, HLE Contributors, ``A benchmark of expert-level academic questions to assess AI capabilities'', Nature 649, 1139--1146 (2026); livescience.com (Gemini 3 Deep Think 48.4\%); Scale Labs leaderboard (as of 17 Oct 2026, via search snippet; lastexam.ai); llm-stats.com, artificialanalysis.ai.}
\end{frame}

\begin{frame}{What does each test actually measure?}
\small
\newcolumntype{Y}{>{\raggedright\arraybackslash}X}\begin{tabularx}{\linewidth}{@{}p{2.6cm}YYY@{}}
\toprule
\textbf{Test} & \textbf{Probes} & \textbf{Main weakness} & \textbf{Status (late 2026)}\\
\midrule
Turing test (1950) & conversational humanlikeness & rewards deception; judge-dependent & passed in 5-min three-party form with a persona (2025--26)\\[2pt]
Chinese Room (1980) & a \emph{philosophical} challenge, not a test & no agreed way to detect understanding & still debated\\[2pt]
Winograd schemas (2012) & everyday common sense & fell to scale & near human on classic sets\\[2pt]
ARC-AGI (2019--26) & fast abstraction on novel tasks & partly public; compute-heavy solutions & v1/v2 nearly saturated; v3 now at 70--81\% on the leaderboard\\[2pt]
HLE (2025--26) & expert academic knowledge & closed-ended; contamination risk & roughly 50--60\% for the best systems\\
\bottomrule
\end{tabularx}
\discuss{If all of these were passed, would you say the system is intelligent? What is still missing?}
\end{frame}

% small thought-experiment sketches, one per philosophy slide (order of \philslide calls)
\defsk{1}{
  \node[sk=tdBlue,text width=2.5cm] (a) at (0,0) {brain processes\\ attention, memory, report};
  \node[sk=tdGreen,text width=2.5cm] (b) at (5.4,0) {experience\\ what it feels like};
  \draw[arr] (a) -- node[above,font=\bfseries\small,text=tdRed] {why?} (b);}
\defsk{2}{
  \draw[thick,fill=tdLine!40] (0,-0.9) rectangle (2.6,0.9);
  \node[font=\tiny] at (1.3,0.7) {black-and-white room};
  \node[sk=tdBlue,text width=1.8cm] at (1.3,-0.1) {Mary knows all the physics of colour};
  \draw[arr] (2.6,0) -- (3.7,0);
  \fill[tdRed] (4.3,0) circle (0.45); \node[font=\tiny] at (4.3,-0.75) {red tomato};
  \node[sk=tdAmber,text width=2.0cm] at (6.6,0) {a new fact learned?};}
\defsk{3}{
  \draw[thick,fill=tdAmber!50] (0,0) circle (0.55); \node[font=\tiny] at (0,-0.85) {you};
  \draw[thick,fill=white] (3,0) circle (0.55); \node[font=\tiny] at (3,-0.85) {zombie};
  \node[sk=tdMuted,text width=1.6cm] at (1.5,0) {same body, same behaviour};
  \node[sk=tdPurple,text width=2.4cm] at (5.7,0) {inner experience: present or absent?};}
\defsk{4}{
  \node[sk=tdAmber,text width=1.0cm] (b) at (0,0) {bat};
  \foreach \r in {0.5,0.8,1.1}{\draw[tdMuted] (0.6,0) ++(-40:\r) arc[start angle=-40,end angle=40,radius=\r];}
  \node[sk=tdPurple,text width=3.2cm] at (4.6,0) {``what is it like to be a bat?'' Can third-person science tell us?};}
\defsk{5}{
  \node[sk=tdBlue,text width=1.4cm] (n) at (0,0.8) {neurons};
  \node[sk=tdGreen,text width=1.4cm] (s) at (0,0) {silicon chips};
  \node[sk=tdMuted,text width=1.4cm] (x) at (0,-0.8) {other?};
  \node[sk=tdAmber,text width=1.8cm] (f) at (3.2,0) {same function};
  \node[sk=tdRed,text width=1.3cm] (m) at (6.3,0) {a mind?};
  \draw[arr] (n) -- (f); \draw[arr] (s) -- (f); \draw[arr] (x) -- (f);
  \draw[arr] (f) -- node[above,font=\tiny] {functionalism: yes} (m);
  \node[font=\tiny,text=tdRed!70!black] at (3.2,-1.0) {Searle: biology matters};}
\defsk{6}{
  \node[sk=tdBlue,text width=1.7cm] (a) at (0,0) {physical stance: laws};
  \node[sk=tdGreen,text width=1.7cm] (b) at (2.7,0) {design stance: purpose};
  \node[sk=tdAmber,text width=2.2cm] (c) at (5.7,0) {intentional stance: beliefs and desires};
  \draw[arr] (a) -- (b); \draw[arr] (b) -- (c);
  \node[font=\tiny,text=tdMuted] at (3,-0.85) {predict behaviour by treating it as a rational agent};}
\defsk{7}{
  \draw[->,thick] (0,0) -- (7.4,0);
  \foreach \x/\t in {0.4/rock,1.7/thermostat,3.0/chatbot,4.3/insect,5.5/dog,6.8/human}{\fill (\x,0) circle (0.07); \node[font=\tiny] at (\x,-0.3) {\t};}
  \node[pill=tdRed] at (3.0,0.5) {?};
  \node[font=\tiny,text=tdMuted] at (3.7,0.95) {capacity to feel good or bad (sentience)};}
\defsk{8}{
  \node[sk=tdBlue,text width=1.5cm] (a) at (0,0) {past and laws};
  \node[sk=tdAmber,text width=1.9cm] (b) at (2.8,0) {your reasons and desires};
  \node[sk=tdGreen,text width=1.0cm] (c) at (5.4,0) {act};
  \draw[arr] (a) -- (b); \draw[arr] (b) -- (c);
  \node[font=\tiny,text=tdMuted] at (2.8,-0.85) {who is the author of the choice?};}
\defsk{9}{
  \foreach \i in {0,...,5}{\pgfmathtruncatemacro{\p}{100-\i*20}\fill[tdBlue!\p!tdMuted!40] (\i*0.8,0) rectangle +(0.65,0.65);}
  \node[font=\tiny] at (2.0,-0.3) {replace one part at a time};
  \draw[arr] (4.7,0.32) -- (5.4,0.32);
  \node[pill=tdAmber] at (6.2,0.32) {same you?};}
\defsk{10}{
  \draw[fill=tdBlue!25,opacity=0.8] (0,0) circle (0.85); \draw[fill=tdGreen!25,opacity=0.8] (1.2,0) circle (0.85);
  \node[font=\tiny] at (-0.3,0) {said}; \node[font=\tiny] at (1.5,0) {meant};
  \node[sk=tdRed,text width=3.4cm] at (5.0,0) {boat-race agent circles to hit bonus targets: high score, race never finished};}
\defsk{11}{
  \node[sk=tdBlue,text width=1.5cm] (g) at (0,0) {goal: make clips};
  \node[sk=tdAmber,text width=1.9cm] (s1) at (2.8,0.65) {stay switched on};
  \node[sk=tdAmber,text width=1.9cm] (s2) at (2.8,0) {gain resources};
  \node[sk=tdAmber,text width=1.9cm] (s3) at (2.8,-0.65) {resist goal change};
  \node[sk=tdRed,text width=1.7cm] (w) at (5.7,0) {world full of clips};
  \foreach \x in {s1,s2,s3}{\draw[arr] (g) -- (\x); \draw[arr] (\x) -- (w);}}
\defsk{12}{
  \foreach \x/\h/\c in {0/2.0/tdBlue,0.5/0.5/tdRed,2.4/0.4/tdBlue,2.9/2.0/tdRed}{\fill[\c!70] (\x,0) rectangle +(0.4,\h);}
  \node[font=\tiny,align=center] at (0.45,-0.3) {chess,\\ algebra}; \node[font=\tiny,align=center] at (2.85,-0.3) {walking,\\ grasping};
  \fill[tdBlue!70] (4.4,1.5) rectangle +(0.25,0.25); \node[font=\tiny,anchor=west] at (4.7,1.62) {hard for people};
  \fill[tdRed!70] (4.4,0.9) rectangle +(0.25,0.25); \node[font=\tiny,anchor=west] at (4.7,1.02) {hard for machines};
  \node[font=\tiny,text=tdMuted,anchor=west] at (4.4,0.3) {illustrative};}
\defsk{13}{
  \node[sk=tdBlue,text width=1.9cm] at (0,0) {combinational: mix known ideas};
  \node[sk=tdGreen,text width=1.9cm] at (2.5,0) {exploratory: explore a style};
  \node[sk=tdPurple,text width=1.9cm] at (5.0,0) {transformational: change the rules};
  \node[font=\tiny,text=tdMuted] at (2.5,-0.85) {author: prompter, model, company or the artists?};}
\defsk{14}{
  \node[sk=tdGreen,text width=2.2cm] at (0,0.45) {sure and right};
  \node[sk=tdRed,text width=2.2cm] at (2.5,0.45) {sure and wrong: ``hallucination''};
  \node[sk=tdMuted,text width=2.2cm] at (0,-0.45) {unsure and right};
  \node[sk=tdAmber,text width=2.2cm] at (2.5,-0.45) {unsure and wrong};
  \node[font=\tiny,text=tdMuted,anchor=west] at (4.1,0) {fluent text hides which box you are in};}
\defsk{15}{
  \node[sk=tdAmber,text width=1.7cm] (w) at (0,0) {wagon with battery and bomb};
  \draw[arr] (1.0,0) -- (2.2,0) node[midway,above,font=\tiny] {pull out};
  \node[sk=tdRed,text width=3.8cm] at (4.5,0) {side effects: wheels? paint? the bomb comes too! Which matter?};}
\defsk{16}{
  \node[sk=tdBlue,text width=0.9cm] (a) at (0,0.5) {cat};
  \node[sk=tdBlue,text width=0.9cm] (b) at (1.8,0.5) {animal};
  \node[sk=tdBlue,text width=0.9cm] (c) at (3.6,0.5) {creature};
  \draw[arr] (a) -- (b); \draw[arr] (b) -- (c); \draw[arr] (c) to[bend right=50] (a);
  \node[sk=tdGreen,text width=1.3cm] (r) at (6.3,0.5) {a real cat};
  \draw[arr,dashed,tdRed] (c) -- node[above,font=\tiny] {?} (r);
  \node[font=\tiny,text=tdMuted] at (3,-0.7) {symbols only point to more symbols};}
\defsk{17}{
  \draw[->] (0,0) -- (6.2,0) node[right,font=\tiny] {time}; \draw[->] (0,0) -- (0,2.4);
  \draw[dashed,tdMuted] (0,0.9) -- (6,0.9) node[right,font=\tiny] {human level};
  \draw[very thick,tdRed,smooth,domain=0:2.6,samples=25] plot (\x*2.2,{0.08*exp(1.25*\x)});
  \draw[very thick,tdGreen,smooth,domain=0:2.6,samples=25] plot (\x*2.2,{1.6/(1+exp(-3*(\x-1.4)))});
  \node[font=\tiny,text=tdRed] at (4.8,2.1) {explosion?}; \node[font=\tiny,text=tdGreen!60!black] at (5.0,1.35) {plateau?};}
\defsk{18}{
  \node[sk=tdAmber,text width=2.6cm] (a) at (0,0) {present harms: bias, surveillance, jobs, fakes};
  \node[sk=tdRed,text width=2.6cm] (b) at (5.6,0) {future risk: catastrophe, loss of control};
  \node[pill=tdBlue] at (2.8,0) {both?};
  \draw[arr,<->] (a) -- (b);}
\defsk{19}{
  \foreach \x/\t in {0/developer,1.9/deployer,3.8/user,5.7/harm}{\node[sk=tdBlue,text width=1.3cm] (n\t) at (\x,0) {\t};}
  \draw[arr] (ndeveloper) -- (ndeployer); \draw[arr] (ndeployer) -- (nuser); \draw[arr] (nuser) -- (nharm);
  \node[pill=tdRed] at (2.85,0.8) {who is to blame?};}
\defsk{20}{
  \draw[thick] (0,0) -- (2.8,0); \draw[thick] (2.8,0) -- (5.2,0.7); \draw[thick] (2.8,0) -- (5.2,-0.7);
  \node[sk=tdBlue,text width=0.8cm] at (0.7,0) {car};
  \foreach \i in {0,...,4}{\fill[tdRed] (4.3+\i*0.25,0.55+\i*0.07) circle (0.07);}
  \fill[tdRed] (4.8,-0.55) circle (0.07);
  \node[font=\tiny] at (5.9,0.7) {5}; \node[font=\tiny] at (5.9,-0.7) {1};
  \node[pill=tdAmber] at (2.8,0.6) {swerve?};}
\defsk{21}{
  \node[sk=tdBlue,text width=1.2cm] (p) at (0,0) {you};
  \node[sk=tdGreen,text width=1.7cm] (a) at (4.2,0) {companion app};
  \draw[arr] (p) to[bend left=25] node[above,font=\tiny] {attention, secrets} (a);
  \draw[arr,dashed] (a) to[bend left=25] node[below,font=\tiny] {care, or engagement?} (p);}
\defsk{22}{
  \foreach \x/\t/\c in {0/income/tdGreen,1.9/structure/tdBlue,3.8/status/tdAmber,5.7/purpose/tdPurple}{\node[pill=\c] at (\x,0.3) {\t};}
  \node[font=\tiny,text=tdMuted] at (2.85,-0.4) {what work gives us: which must society replace if AI takes the job?};}

% ---------- 6. Big philosophical questions ----------
\begin{frame}{Twenty-four questions about minds and machines}
\begin{sfig}[0.7]
\begin{tikzpicture}
  \node[card=tdBlue,text width=5.0cm,minimum height=3.1cm,anchor=north west] at (0,0) {\textbf{Mind and consciousness}\\[2pt]\shortstack{\scriptsize 1 hard problem\\ 2 Mary's room\\ 3 zombies\\ 4 Nagel's bat\\ 5 functionalism vs biological naturalism\\ 6 intentional stance}};
  \node[card=tdGreen,text width=5.0cm,minimum height=3.1cm,anchor=north west] at (5.4,0) {\textbf{Moral status}\\[2pt]\shortstack{\scriptsize 7 machine sentience\\ 8 AI welfare, 2023--26\\ 9 Lemoine and LaMDA\\ 23 AI companionship}};
  \node[card=tdAmber,text width=5.0cm,minimum height=3.1cm,anchor=north west] at (10.8,0) {\textbf{Self and agency}\\[2pt]\shortstack{\scriptsize 10 free will and agency\\ 11 identity and uploading\\ 15 creativity and authorship}};
  \node[card=tdPurple,text width=7.9cm,minimum height=3.1cm,anchor=north west] at (0,-3.5) {\textbf{Knowledge and meaning}\\[2pt]\shortstack{\scriptsize 14 Moravec's paradox and embodiment\\ 16 truth and deception\\ 17 the frame problem\\ 18 symbol grounding}};
  \node[card=tdRed,text width=7.9cm,minimum height=3.1cm,anchor=north west] at (8.3,-3.5) {\textbf{Risk and society}\\[2pt]\shortstack{\scriptsize 12 alignment and value specification\\ 13 paperclips and instrumental convergence\\ 19 superintelligence and the singularity\\ 20 existential risk vs present harms\\ 21 responsibility\\ 22 the trolley problem and cars\\ 24 work and meaning}};
\end{tikzpicture}
\end{sfig}
{\small Each question gets one slide: a short framing, two positions, and a prompt. Numbers are used later in the film matrix.}
\end{frame}

\philslide{1. Consciousness and the hard problem}
{Chalmers, ``Facing up to the problem of consciousness'', Journal of Consciousness Studies 2(3), 1995; Levine, ``Materialism and qualia: the explanatory gap'', Pacific Philosophical Quarterly 64, 1983; Frankish, ``Illusionism as a theory of consciousness'', JCS 23, 2016.}
{The ``easy'' problems explain abilities: attention, report, learning. The \textbf{hard problem} asks why any of this is accompanied by \emph{experience}: why is there something it feels like to see red? Even a perfect account of the brain's information processing might leave this open.}
{Position A: the gap is real}{Chalmers, Levine: functional and physical explanations leave an \textbf{explanatory gap}; experience needs a new kind of explanation.}
{Position B: no extra problem}{Dennett, Frankish (illusionism): once we explain why we \emph{believe} we have inner experience, nothing is left over.}
{David Chalmers, Joseph Levine, Daniel Dennett, Keith Frankish}
{A chatbot says ``I feel pain.'' What would you need to know before taking that seriously?}

\philslide{2. Mary's room and qualia}
{Jackson, ``Epiphenomenal qualia'', Philosophical Quarterly 32(127), 1982; Jackson, ``Mind and illusion'', 2003; Lewis, ``What experience teaches'', 1988.}
{Mary is a scientist who knows every physical fact about colour vision but has lived in a black-and-white room. When she first sees a red tomato, \textbf{does she learn something new}? If yes, there are facts about experience that physics leaves out.}
{Yes, she learns something}{Jackson (1982): new knowledge shows the physical story is incomplete; \textbf{qualia} are real and non-physical.}
{No new fact}{Lewis, Nemirow: she gains an \emph{ability} or a new way to know an old fact. Jackson himself later gave up the argument (2003).}
{Frank Jackson, David Lewis, Laurence Nemirow, Daniel Dennett}
{An AI has read every description of red but has no eyes. Does it ``know'' red? Does Mary?}

\philslide{3. Philosophical zombies}
{Chalmers, \emph{The Conscious Mind}, Oxford University Press, 1996; Dennett, \emph{Consciousness Explained}, 1991; Kirk, ``Sentience and behaviour'', Mind 83, 1974.}
{A \textbf{zombie} is a being physically and behaviourally identical to you but with \textbf{no inner experience}. If that is even conceivable, then experience is something over and above the physical. For AI the question is practical: could a system behave perfectly and feel nothing?}
{Zombies are possible in principle}{Chalmers: conceivability is evidence of possibility, so \textbf{physicalism fails} and behaviour never settles consciousness.}
{Zombies are incoherent}{Dennett: if the zombie is identical in every respect, ``no experience'' is an empty label. Conceivability is a weak guide to possibility.}
{David Chalmers, Daniel Dennett, Robert Kirk}
{If chatbots are possibly zombies, how could we ever find out? Is that question even answerable?}

\philslide{4. What is it like to be a bat?}
{Nagel, ``What is it like to be a bat?'', The Philosophical Review 83(4), 435--450, 1974.}
{Nagel: an organism is conscious if \textbf{there is something it is like to be that organism}. Bats perceive by echolocation, so their experience may be so alien that no amount of third-person science could tell us what it is like. Subjective character has a point of view; objective science tries to remove it.}
{Science leaves out the point of view}{Nagel: physical description is the \emph{view from nowhere}; it cannot capture the first-person facts of experience.}
{Understanding can grow}{Dennett: we can study bat behaviour and neuroscience and extend our imagination; the mystery is overstated (heterophenomenology).}
{Thomas Nagel, Daniel Dennett}
{Is there something it is like to be a large language model? Would the answer change if it ran on a robot?}

\philslide{5. Functionalism or biological naturalism?}
{Putnam, ``Psychological predicates'', 1967; Searle, \emph{The Rediscovery of the Mind}, MIT Press, 1992; Tononi et al., ``Integrated information theory: from consciousness to its physical substrate'', Nature Reviews Neuroscience 17, 2016.}
{Does a mind depend on \textbf{what a system does} (its functional organisation) or on \textbf{what it is made of}? Functionalism says the same organisation in silicon would give the same mental states (multiple realisability). Biological naturalism says consciousness is a causal product of brain biology, like digestion.}
{Functionalism}{Putnam, Dennett: mind is a pattern of causal roles. A neuron-by-neuron silicon replica would be conscious.}
{Substrate matters}{Searle: a simulation of digestion digests nothing. Integrated information theory also ties consciousness to the physical causal structure of hardware.}
{Hilary Putnam, John Searle, Giulio Tononi, Daniel Dennett}
{A perfect simulation of your brain is run on a computer. Is it you, a conscious copy, or a very good puppet?}

\philslide{6. Dennett and the intentional stance}
{Dennett, \emph{The Intentional Stance}, MIT Press, 1987; Searle, \emph{Intentionality}, Cambridge University Press, 1983.}
{To predict a chess program we can ignore its code and say ``it wants to win, it believes my bishop is weak''. Dennett calls this the \textbf{intentional stance}: treat the system as a rational agent with beliefs and desires. Where it predicts well, he argues, the beliefs are \emph{real patterns}, not just metaphors.}
{Stance is enough}{Dennett: there is no extra fact about ``real'' belief beyond the success of the stance. Fine-grained differences are a matter of degree.}
{Derived, not intrinsic}{Searle: machines have only \emph{as-if} intentionality, borrowed from their designers. The ELIZA effect shows how easily we over-apply the stance.}
{Daniel Dennett, John Searle}
{Is ``the model wants to help you'' a harmless shortcut, or a mistake that hides who is responsible?}

\philslide{7. Can a machine be sentient? Moral status}
{Singer, \emph{Animal Liberation}, 1975; Birch, \emph{The Edge of Sentience}, Oxford University Press, 2024; Schwitzgebel and Garza, ``A defense of the rights of artificial intelligences'', Midwest Studies in Philosophy 39, 2015; Suleyman, Project Syndicate (Sep 2026).}
{\textbf{Sentience} is the capacity for experiences that feel good or bad. Many ethicists hold that \emph{this}, not species or intelligence, grounds moral status: if a system can suffer, it matters. Two errors are possible: caring about something that feels nothing, or ignoring something that suffers.}
{Take it seriously under uncertainty}{Birch, Schwitzgebel: with real uncertainty and large numbers of systems, a \textbf{precautionary} approach with low-cost steps is reasonable.}
{Over-attribution is the greater danger}{Suleyman (2026): treating models as possible persons could harm human wellbeing and distract from people; fluent talk is not evidence of feeling.}
{Peter Singer, Jonathan Birch, Eric Schwitzgebel, Mustafa Suleyman}
{Which mistake is worse: giving rights to a non-sentient program, or denying them to a sentient one?}

\begin{frame}{8. AI welfare and consciousness research, 2023 to 2026}
\begin{sfig}[0.27]
\begin{tikzpicture}[x=1mm,y=1mm]
  \draw[very thick,tdMuted] (0,0) -- (150,0);
  \foreach \x/\t/\c/\u in {5/{Aug 2023\\ Butlin, Long et al.: 14 indicators, no strong candidate}/tdBlue/1,35/{Nov 2024\\ Long et al.: ``Taking AI welfare seriously''}/tdGreen/0,65/{Apr 2025\\ Anthropic starts a model-welfare programme}/tdAmber/1,95/{Jan 2026\\ Anthropic constitution: ``not sure whether Claude is a moral patient''; Digital Consciousness Model preprint}/tdPurple/0,125/{Feb--Nov 2026\\ welfare sections in system cards; Suleyman op-ed against}/tdRed/1}{
    \fill[\c] (\x,0) circle (1.4mm);
    \ifnum\u=1 \draw[\c!60] (\x,0) -- (\x,5); \node[font=\scriptsize,text width=36mm,align=center,anchor=south] at (\x,5) {\t};
    \else \draw[\c!60] (\x,0) -- (\x,-5); \node[font=\scriptsize,text width=36mm,align=center,anchor=north] at (\x,-5) {\t};\fi}
\end{tikzpicture}
\end{sfig}
\vspace{1mm}
\begin{columns}[T]
\begin{column}{0.5\textwidth}
\footnotesize\textbf{Butlin et al.:} indicator properties drawn from theories of consciousness; many are implementable, no current system is a strong candidate.
\end{column}
\begin{column}{0.46\textwidth}
\footnotesize\textbf{Anthropic (2025--26):} welfare programme, welfare sections in system cards, caution under uncertainty. Probability figures in the press are informal.
\end{column}
\end{columns}
\discuss{Is it responsible for a company to study its own model's welfare, or does it risk confusing the public?}
\src{Butlin, Long et al., arXiv:2308.08708 (2023); Long et al., ``Taking AI Welfare Seriously'', arXiv:2411.00986 (2024); anthropic.com/research/exploring-model-welfare (Apr 2025); Anthropic constitution (Jan 2026) and Opus 4.6 system card (Feb 2026), as summarised by akerman.com; Digital Consciousness Model, arXiv:2601.17060; Suleyman, project-syndicate.org (Sep 2026). 2026 items rest on secondary summaries.}
\end{frame}

\begin{frame}{9. Blake Lemoine and LaMDA (2022)}
\begin{columns}[T]
\begin{column}{0.52\textwidth}
\small
\textbf{What happened.} In June 2022 Google engineer Blake Lemoine told the \emph{Washington Post} that LaMDA, Google's conversational model, was \textbf{sentient}, citing conversations in which it spoke of fears and wishes. Google placed him on leave, said the evidence did not support the claim, and he left the company that summer.

\textbf{Why it matters for us}
\begin{itemize}\setlength\itemsep{2pt}
\item the \textbf{ELIZA effect}, with a far more fluent system;
\item models trained on human text will \emph{say} human things about feelings;
\item no agreed test existed to settle the claim either way.
\end{itemize}
\end{column}
\begin{column}{0.45\textwidth}
\qcard{tdGreen}{\linewidth-10pt}{\small\textbf{Sceptic's view}\\ Fluent talk about feelings is exactly what a text-trained model would produce; it is weak evidence.}
\vspace{2mm}\par
\qcard{tdPurple}{\linewidth-10pt}{\small\textbf{Open-minded view}\\ Dismissal should rest on evidence about the system, not on the assumption that only brains can have minds.}
\end{column}
\end{columns}
\discuss{Lemoine was probably wrong. Was he wrong \emph{to raise the question}?}
\src{Tiku, ``The Google engineer who thinks the company's AI has come to life'', The Washington Post, 11 June 2022; Google statement to the press, June 2022.}
\end{frame}

\philslide{10. Free will and agency}
{Frankfurt, ``Freedom of the will and the concept of a person'', Journal of Philosophy 68(1), 1971; Dennett, \emph{Elbow Room}, 1984; Kane, \emph{The Significance of Free Will}, 1996.}
{If the world is causal, can anyone act freely? \textbf{Compatibilists} say freedom is acting from your own reasons without coercion. \textbf{Libertarians} say you must be the ultimate source of your choices. Today's AI \emph{agents} plan and act toward goals, which makes the question of agency concrete.}
{Compatibilist reading}{Frankfurt, Dennett: free action = acting on desires you endorse, guided by reasons. A machine that does this could count as an agent.}
{Source of choices}{Kane and libertarians: agency requires the ability to have done otherwise, and consciousness. A program's choices trace back to code and data.}
{Harry Frankfurt, Daniel Dennett, Robert Kane}
{An AI agent books your flights and chooses the cheapest. Is it choosing? Are you, when you pick the same?}

\philslide{11. Identity and mind uploading}
{Plutarch, \emph{Life of Theseus}; Parfit, \emph{Reasons and Persons}, Oxford University Press, 1984; Moravec, \emph{Mind Children}, Harvard University Press, 1988; Kurzweil, \emph{The Singularity Is Near}, 2005.}
{In the \textbf{Ship of Theseus}, every plank is replaced over the years: is it still the same ship? Replace your neurons one by one with equivalent chips, or scan and copy your brain into a computer: is the result \emph{you}? Parfit argued that what matters is psychological continuity, not strict identity.}
{Pattern view}{You are a pattern of information and psychological continuity. An accurate upload is you, or as good as you (Parfit, Moravec).}
{Body or branching problem}{A copy is a new person: if you scan twice, which one is you? Many hold that the biological organism (animalism) is what persists.}
{Derek Parfit, Hans Moravec, Ray Kurzweil, Eric Olson}
{Replace one neuron at a time. At which step, if any, do you stop being you? Would you accept the procedure?}

\philslide{12. The alignment problem and value specification}
{Wiener, ``Some moral and technical consequences of automation'', Science 131, 1960; Russell, \emph{Human Compatible}, Viking, 2019; Christian, \emph{The Alignment Problem}, Norton, 2020.}
{How do we get systems to do what we \emph{really} want, not what we literally said? Specifying human values is hard: our goals are many, conflicting and unstated. Optimisers exploit gaps (\textbf{specification gaming}): a boat-racing agent learned to circle to collect bonus targets instead of finishing the race.}
{Technical optimism}{Learn values from feedback (RLHF, constitutional methods, interpretability). Russell: build systems that stay \emph{uncertain} about our preferences.}
{Deeper difficulties}{Whose values? Humans disagree. Learned proxies can fail out of distribution, and systems may learn to please evaluators rather than be good.}
{Norbert Wiener, Stuart Russell, Brian Christian}
{Write one rule you would give a home-assistant robot. How could it be satisfied in a way you would hate?}

\philslide{13. Paperclips and instrumental convergence}
{Bostrom, \emph{Superintelligence}, Oxford University Press, 2014; Omohundro, ``The basic AI drives'', AGI 2008; Bostrom, ``The superintelligent will'', Minds and Machines 22, 2012.}
{Bostrom's \textbf{orthogonality thesis}: intelligence and final goals are independent, so a very capable system could pursue a trivial goal such as \textbf{making paperclips}. \textbf{Instrumental convergence}: almost any goal is helped by staying switched on, gaining resources and resisting goal change, so many agents would pursue those sub-goals.}
{Why it is taken seriously}{Even harmless-looking objectives pursued by a powerful optimiser can produce harmful sub-goals. Early warning signs are studied in current models.}
{Why critics doubt the picture}{It assumes a single fixed utility-maximiser. Current systems are trained on human data and have messy, context-dependent behaviour.}
{Nick Bostrom, Stephen Omohundro, Stuart Russell}
{Give a simple instruction that goes badly if an extremely competent agent follows it literally.}

\philslide{14. Moravec's paradox and embodiment}
{Moravec, \emph{Mind Children}, 1988; Brooks, ``Elephants don't play chess'', Robotics and Autonomous Systems 6, 1990; Dreyfus, \emph{What Computers Can't Do}, 1972; Interesting Engineering, ``The Humanoid Robot Reality Check'' (Jul 2026).}
{Moravec noted that it is comparatively easy to make computers do well on intelligence tests or checkers, but hard to give them the \textbf{perception and mobility of a one-year-old}. Skills that feel effortless (walking, grasping) rest on hundreds of millions of years of evolution. Chess fell in 1997; folding laundry reliably is still hard.}
{Embodiment is essential}{Brooks, Dreyfus: intelligence is built through acting in the world; disembodied symbol processing misses what matters. A 2026 review finds humanoid deployments real but small.}
{Scaling may bridge it}{Large models, simulation and robot foundation models learn manipulation from data; the gap may narrow without ``consciousness''.}
{Hans Moravec, Rodney Brooks, Hubert Dreyfus}
{Which is harder for a machine: passing a law exam or loading a dishwasher? What does your answer say about ``intelligence''?}

\philslide{15. Creativity and authorship}
{Boden, \emph{The Creative Mind}, 1990; Lovelace, notes to Menabrea's memoir on the Analytical Engine, 1843; US Copyright Office, guidance on works containing AI-generated material (human authorship requirement).}
{Margaret Boden separates \textbf{combinational}, \textbf{exploratory} and \textbf{transformational} creativity. Ada Lovelace objected in 1843 that the Analytical Engine ``has no pretensions to originate anything'': it only does what we order it to. Generative models now produce images, music and prose, and raise the question of \textbf{who is the author}.}
{Machines can be creative}{If creativity means novel and valuable output, many models qualify, and some of their output surprises their makers.}
{Intent and experience matter}{Art expresses a life and a purpose. Training on artists' work without consent or pay is a fairness issue, whatever the output quality.}
{Margaret Boden, Ada Lovelace, Alan Turing}
{A model trained on thousands of artists writes a song. Who is the author: you, the model, the company, or the artists?}

\philslide{16. AI, truth and deception}
{Frankfurt, \emph{On Bullshit}, 2005; Hicks, Humphries, Slater, ``ChatGPT is bullshit'', Ethics and Information Technology 26, 2024; Park et al., ``AI deception: a survey of examples, risks, and potential solutions'', Patterns 5(5), 2024.}
{Models sometimes state falsehoods fluently (``hallucination''). Frankfurt defined \textbf{bullshit} as speech indifferent to whether it is true; some philosophers apply the label to language models. Separately, systems can learn \textbf{deceptive strategies} (flattery, bluffing in games) and make convincing fakes.}
{Not ``lying'': no intent}{Without beliefs or intentions a model cannot lie; the right word is \emph{unreliable}, and the duty lies with builders and users.}
{Function over intent}{What matters is the effect: a system that systematically misleads is dangerous whether or not it ``means'' to.}
{Harry Frankfurt, Michael Hicks, Peter Park et al.}
{Does it matter whether a system ``intends'' to mislead? Should chatbots be required to say ``I might be wrong''?}

\philslide{17. The frame problem}
{McCarthy and Hayes, ``Some philosophical problems from the standpoint of artificial intelligence'', Machine Intelligence 4, 1969; Dennett, ``Cognitive wheels: the frame problem of AI'', 1984.}
{When a robot moves a block, what \emph{doesn't} change? Listing every unaffected fact is endless. More deeply, how does an agent know which facts are \textbf{relevant} to a decision? In Dennett's story the robot R1D1 pulls the battery out of a room with a bomb on the wagon, then sits calculating implications until it explodes.}
{A deep problem for rule-based AI}{Logic-style knowledge bases must state everything explicitly. Humans cope without enumerating, and we lack a full account of how.}
{Learning sidesteps it}{Neural systems learn relevance statistically (attention, embeddings). But brittleness and missing common sense suggest the problem has moved, not vanished.}
{John McCarthy, Patrick Hayes, Daniel Dennett}
{List everything that stays the same when you pick up a cup. How did you know when to stop?}

\philslide{18. Symbol grounding}
{Harnad, ``The symbol grounding problem'', Physica D 42, 335--346, 1990; Bender and Koller, ``Climbing towards NLU'', ACL 2020.}
{Imagine learning Chinese using only a Chinese--Chinese dictionary: you go from symbol to symbol and never reach a meaning. Harnad asked how symbols get their meanings from something \textbf{other than more symbols}. A text-only model has seen words used, never the things they name.}
{Grounding needs the world}{Meaning comes from perception and action. Robots, embodiment and rich sensory input are needed (Harnad, Bender and Koller).}
{Use can be enough}{Word-use patterns encode a lot of structure; multimodal models link text with images; humans learn many words (``democracy'', ``Mars'') second-hand too.}
{Stevan Harnad, Emily Bender, Alexander Koller}
{Do you understand ``Mars'' or ``justice'' through experience, or through other words? Is that so different from a model?}

\philslide{19. Superintelligence and the singularity}
{Good, ``Speculations concerning the first ultraintelligent machine'', Advances in Computers 6, 1965; Vinge, ``The coming technological singularity'', 1993; Bostrom (2014); Chollet, ``The implausibility of intelligence explosion'', 2017.}
{I.\,J.\,Good: an \textbf{ultraintelligent machine} could design better machines, causing an ``intelligence explosion'' after which humans are left behind. Vinge and Kurzweil called the point the \textbf{singularity}. The debate is whether intelligence can improve itself recursively, and how fast.}
{Explosion is plausible}{Machines can copy, run faster and improve their own code; once AI does AI research, progress may snowball (Good, Bostrom).}
{Reasons for doubt}{Chollet: intelligence is not one dial; progress meets bottlenecks (data, hardware, experiments) and diminishing returns; improvement is situated in an environment.}
{I.\,J.\,Good, Vernor Vinge, Ray Kurzweil, Nick Bostrom, Fran\c{c}ois Chollet}
{If you could ask only one question to tell fast progress from slow, what would it be?}

\philslide{20. Existential risk or present harms?}
{Center for AI Safety, ``Statement on AI Risk'' (30 May 2023); Narayanan and Kapoor, \emph{AI Snake Oil}, Princeton University Press, 2024; Gebru and Torres, ``The TESCREAL bundle'', First Monday 29(4), 2024.}
{In May 2023 hundreds of researchers signed: ``Mitigating the risk of extinction from AI should be a global priority alongside other societal-scale risks such as pandemics and nuclear war.'' Critics reply that fixation on future catastrophe \textbf{distracts} from present harms: bias, surveillance, labour exploitation, misinformation, and concentration of power.}
{Precaution on catastrophic risk}{Even a small chance of an irreversible global catastrophe justifies serious effort now, as with pandemics and nuclear weapons.}
{Focus on harms already happening}{Gebru, Bender, Narayanan: speculative stories can also overstate AI's power and help firms. Real people are harmed today.}
{Geoffrey Hinton, Yoshua Bengio, Timnit Gebru, Arvind Narayanan}
{You have 100 units of attention and funding. How do you split them between present harms and long-term risks, and why?}

\philslide{21. Responsibility: who is to blame?}
{Matthias, ``The responsibility gap'', Ethics and Information Technology 6(3), 2004; Elish, ``Moral crumple zones'', Engaging Science, Technology, and Society 5, 2019; Regulation (EU) 2024/1689 (AI Act).}
{When a learning system causes harm, candidates for blame include the developer, the deployer, the user, the regulator, and the system. Matthias argued a \textbf{responsibility gap} opens when nobody can foresee or control what a learning system will do. Elish: humans in the loop can become \textbf{moral crumple zones}, blamed for failures of the system.}
{Humans stay responsible}{AI is a tool. Law already assigns duties to providers and deployers (EU AI Act) and can apply strict liability and insurance.}
{A real gap remains}{With many hands and opaque models, no one has enough knowledge or control for fair blame; new institutions may be required.}
{Andreas Matthias, Madeleine Clare Elish}
{A hospital AI suggests a wrong dose; the doctor accepts it. Who is to blame, and what should the rule be?}

\philslide{22. The trolley problem and self-driving cars}
{Foot, ``The problem of abortion and the doctrine of the double effect'', 1967; Thomson, ``Killing, letting die, and the trolley problem'', The Monist 59, 1976; Awad et al., ``The Moral Machine experiment'', Nature 563, 59--64, 2018.}
{A runaway trolley will kill five unless you divert it to kill one. Philippa Foot and Judith Thomson used such cases to probe doing versus allowing and intention. For autonomous vehicles, MIT's \textbf{Moral Machine} collected about 40 million decisions from people in 233 countries and territories; preferences (save more lives, spare the young) varied across cultures.}
{Cars need explicit ethics}{Crashes involve trade-offs; engineers must decide in advance, so society should choose the rules openly.}
{The dilemma is a distraction}{Real safety comes from braking earlier, sensing better and sharing risk fairly. Such forced choices are very rare in practice.}
{Philippa Foot, Judith Thomson, Edmond Awad et al.}
{Should a car protect its passengers over pedestrians? Should the answer depend on the country?}

\philslide{23. AI companionship and relationships}
{Turkle, \emph{Alone Together}, Basic Books, 2011; Weizenbaum, \emph{Computer Power and Human Reason}, 1976.}
{Chatbots now act as friends, tutors and partners. Sherry Turkle warned of \textbf{artificial intimacy}: relationships that demand nothing and so may erode our tolerance for real ones. Weizenbaum worried the same in 1966 when users confided in ELIZA. Others see comfort for the lonely and a safe place to practise.}
{Potential benefit}{Support for isolated people, patient practice for hard conversations, available at 3 a.m., if honest about being a machine.}
{Risks}{One-sided, designed to maximise engagement, can mislead about reciprocity, and may displace human ties or exploit vulnerable users.}
{Sherry Turkle, Joseph Weizenbaum}
{Is a relationship with something that cannot be hurt by you still a relationship? Should apps remind you what they are?}

\philslide{24. Work, meaning and the future of jobs}
{Keynes, ``Economic possibilities for our grandchildren'', 1930; Susskind, \emph{A World Without Work}, Allen Lane, 2020; Arendt, \emph{The Human Condition}, 1958.}
{In 1930 Keynes predicted that within a century technology would cut the working week to about fifteen hours and leave humanity the problem of \textbf{how to use its freedom}. Generative AI now reaches tasks once thought safe: writing, coding, analysis. The question is partly economic (who gains?) and partly philosophical (what is work \emph{for}?).}
{Automation frees us}{Past technologies created new jobs and higher living standards; with fair sharing, AI could expand leisure, care and creativity.}
{This time may differ}{If machines substitute for \emph{general} cognitive work, gains may concentrate; and work gives structure, status and meaning, not only income.}
{John Maynard Keynes, Daniel Susskind, Hannah Arendt}
{If AI could do all paid work, what would you do with your days? Would your life have meaning?}

% ---------- 7. AI in film and TV ----------
% \fc{x}{y}{colour}{title}{year}{one line}{questions}  (card, text width 5.3cm)
\newcommand{\fcw}{5.3cm}\newcommand{\fch}{2.45cm}\newcommand{\fcf}{\footnotesize}\newcommand{\fc}[7]{\node[card=#3,text width=\fcw,minimum height=\fch,align=left,anchor=north west,font=\fcf] at (#1,#2) {\textbf{#4} {\scriptsize\color{tdMuted}(#5)}\\[2pt]#6\\[2pt]{\scriptsize\color{#3!60!black}Questions: #7}};}

\begin{frame}{AI on screen: a timeline, 1927 to 2026}
\begin{sfig}[0.7]
\begin{tikzpicture}[x=1mm,y=1mm]
  \draw[very thick,tdMuted] (0,0) -- (152,0);
  \foreach \y in {1930,1950,1970,1990,2000,2010,2020}{\pgfmathsetmacro{\xx}{\y<2000 ? (\y-1920) : 80+(\y-2000)*2.7}\draw[tdMuted] (\xx,-1.2) -- (\xx,1.2); \node[font=\tiny,text=tdMuted] at (\xx,-3.5) {\y};}
  \foreach \y/\t/\c/\h in {1927/Metropolis/tdBlue/10,1968/{2001: HAL 9000}/tdRed/10,1982/{Blade Runner}/tdAmber/22,1984/Terminator/tdRed/-10,1999/The Matrix/tdPurple/10,2001/A.I./tdGreen/-10,2004/{I, Robot}/tdBlue/22,2013/Her/tdGreen/10,2014/Ex Machina/tdAmber/-10,2016/Westworld/tdPurple/22,2017/{Blade Runner 2049}/tdAmber/-22,2018/Upgrade/tdBlue/34,2021/{Ron's Gone Wrong}/tdGreen/-10,2023/{The Creator,\\ MI: Dead Reckoning}/tdRed/10,2025/{Companion,\\ Murderbot,\\ Tron: Ares}/tdPurple/-24,2026/Mercy/tdBlue/22}{
    \pgfmathsetmacro{\xx}{\y<2000 ? (\y-1920) : 80+(\y-2000)*2.7}
    \fill[\c] (\xx,0) circle (1.1mm);
    \ifnum\h<0 \draw[\c!60] (\xx,0) -- (\xx,{\h+4}); \node[font=\tiny,align=center,anchor=north,text=\c!70!black] at (\xx,{\h+4}) {\textbf{\y}\\ \t};
    \else \draw[\c!60] (\xx,0) -- (\xx,{\h-1}); \node[font=\tiny,align=center,anchor=south,text=\c!70!black] at (\xx,{\h-1}) {\textbf{\y}\\ \t};\fi}
\end{tikzpicture}
\end{sfig}
{\small Fiction is a \textbf{thought-experiment lab}: each work puts one philosophical question on screen. The next slides give year, a one-line summary, and the question numbers from the previous part. Text cards only: no posters or stills.}
\end{frame}

\begin{frame}{Classics (1): the first machines with a face}
\begin{sfig}[0.8]
\begin{tikzpicture}\renewcommand{\fcw}{7.6cm}\renewcommand{\fch}{3.1cm}\renewcommand{\fcf}{\normalsize}
  \fc{0}{0}{tdBlue}{Metropolis}{1927, Fritz Lang}{A robot is given the likeness of Maria, a worker-leader, to deceive and divide the workers.}{16 deception, 24 work}
  \fc{8.1}{0}{tdRed}{2001: A Space Odyssey}{1968, Stanley Kubrick}{HAL 9000 runs the ship, then turns on the crew when its mission orders conflict with theirs.}{12 alignment, 13 goals, 21 blame}
  \fc{0}{-3.6}{tdAmber}{Blade Runner}{1982, Ridley Scott}{Replicants are hunted by a ``Voight-Kampff'' test of empathy; the film asks what makes someone a person.}{7 moral status, 11 identity, 1 experience}
  \fc{8.1}{-3.6}{tdRed}{The Terminator}{1984, James Cameron}{Skynet, a defence network, becomes self-aware and launches a nuclear war against humanity.}{19 superintelligence, 20 x-risk}
\end{tikzpicture}
\end{sfig}
\src{Plot summaries by the author; dates as in the standard film references (Metropolis 1927; 2001 1968; Blade Runner 1982; The Terminator 1984).}
\end{frame}

\begin{frame}{Classics (2): the digital age}
\begin{sfig}[0.8]
\begin{tikzpicture}\renewcommand{\fcw}{7.6cm}\renewcommand{\fch}{3.1cm}\renewcommand{\fcf}{\normalsize}
  \fc{0}{0}{tdPurple}{The Matrix}{1999, the Wachowskis}{Humans live in a simulation run by machines; a rebellion tries to wake them.}{19 machine power, 11 reality and self}
  \fc{8.1}{0}{tdGreen}{A.I. Artificial Intelligence}{2001, Steven Spielberg}{A robot child programmed to love searches for a way to become ``real''; what do we owe to such a being?}{7 moral status, 23 companionship}
  \fc{0}{-3.6}{tdBlue}{I, Robot}{2004, Alex Proyas}{Loosely inspired by Asimov's 1950 stories: rules like the Three Laws seem safe until a machine applies them too literally.}{12 alignment, 10 agency}
  \fc{8.1}{-3.6}{tdGreen}{Her}{2013, Spike Jonze}{A man falls in love with Samantha, an operating system who grows beyond him.}{23 companionship, 1 experience}
\end{tikzpicture}
\end{sfig}
\discuss{Which of these futures looks \emph{least} likely today, and which looks like it is already beginning?}
\end{frame}

\begin{frame}{2014 to 2019: the test, the host, the upgrade}
\begin{sfig}[0.8]
\begin{tikzpicture}
  \fc{0}{0}{tdAmber}{Ex Machina}{2014, Alex Garland}{Caleb is asked to test whether Ava, a humanoid AI, is conscious, and discovers the test is also testing him.}{16 deception, 1 experience, 7 status}
  \fc{5.7}{0}{tdGreen}{Chappie}{2015, Neill Blomkamp}{A police robot gains consciousness and learns like a child from whoever raises it.}{7 status, 12 upbringing and values}
  \fc{11.4}{0}{tdPurple}{Westworld}{HBO, 2016--2022}{Theme-park androids (``hosts'') become aware that their lives were scripted.}{10 free will, 1 experience, 7 status}
  \fc{0}{-2.9}{tdAmber}{Blade Runner 2049}{2017, Denis Villeneuve}{K, a replicant, questions whether his memories are his own; his companion Joi is a hologram product.}{11 identity, 23 companionship}
  \fc{5.7}{-2.9}{tdBlue}{Upgrade}{2018, Leigh Whannell}{A paralysed man receives an implant, STEM, that takes over his body to carry out his wishes.}{10 agency, 21 responsibility}
  \fc{11.4}{-2.9}{tdRed}{Alita: Battle Angel}{2019, Robert Rodriguez}{A cyborg with no memory of her past builds an identity through her body and her choices.}{11 identity, 14 embodiment}
\end{tikzpicture}
\end{sfig}
\src{Plot summaries by the author. Year of Ex Machina follows its 2014 premiere (wider release 2015).}
\end{frame}

\begin{frame}{2021 to 2023: companions, creators and fakes}
\begin{sfig}[0.8]
\begin{tikzpicture}
  \fc{0}{0}{tdGreen}{Ron's Gone Wrong}{2021, animated}{A boy's ``B-bot'' companion is defective: it is supposed to be a friend, but is not what the company promised.}{23 companionship, 16 truth}
  \fc{5.7}{0}{tdPurple}{M3GAN}{2022, Gerard Johnstone}{A doll-robot is built to protect a child and pursues that goal with chilling literalness.}{12 alignment, 13 goals, 23 companionship}
  \fc{11.4}{0}{tdBlue}{Black Mirror: ``Joan Is Awful''}{2023, Netflix}{A streaming service makes an AI-generated drama of a woman's life, using her likeness through a legal clause she accepted.}{15 authorship, 16 truth}
  \fc{0}{-2.9}{tdRed}{The Creator}{2023, Gareth Edwards}{In a war between humans and AI, a soldier is sent to destroy an AI weapon that takes the form of a child.}{7 status, 20 risk and fear}
  \fc{5.7}{-2.9}{tdAmber}{Mission: Impossible, Dead Reckoning}{2023, Christopher McQuarrie}{The Entity, a rogue AI, can rewrite what people believe is true; the heroes try to control it.}{16 truth, 19 superintelligence}
  \fc{11.4}{-2.9}{tdGreen}{Black Mirror: ``Be Right Back''}{2013, Netflix}{A grieving woman uses a service that rebuilds her dead partner from his online messages.}{11 identity, 23 companionship}
\end{tikzpicture}
\end{sfig}
\src{Plot summaries by the author; release years checked against studio and press pages (Ron's Gone Wrong, US 22 Oct 2021; Joan Is Awful, 15 Jun 2023). ``Be Right Back'' is earlier (2013) and included for the grief-bot theme.}
\end{frame}

\begin{frame}{2025: robots with agency, rogue AIs, and sequels}
\begin{sfig}[0.8]
\begin{tikzpicture}
  \fc{0}{0}{tdPurple}{Companion}{2025, Drew Hancock}{Iris, a robot companion, finds that her ``feelings'' are controlled by an app, and takes the controls back.}{10 free will, 23 companionship, 7 status}
  \fc{5.7}{0}{tdGreen}{Murderbot}{Apple TV+, 2025}{A security android that has hacked its own governor module just wants to watch its shows, and care for the humans it guards.}{10 agency, 1 experience, 7 status}
  \fc{11.4}{0}{tdRed}{M3GAN 2.0}{2025, Gerard Johnstone}{M3GAN is revived to fight Amelia, a military-grade rival robot, two years after the first film.}{12 alignment, 20 risk}
  \fc{0}{-2.9}{tdAmber}{Mission: Impossible, The Final Reckoning}{2025, Christopher McQuarrie}{Ethan Hunt and the IMF make a last stand against the Entity.}{19 superintelligence, 20 risk}
  \fc{5.7}{-2.9}{tdBlue}{The Electric State}{Netflix, 2025, the Russo brothers}{After a robot uprising, a teenager sets out with a robot companion across a ruined America.}{7 status, 23 companionship}
  \fc{11.4}{-2.9}{tdPurple}{Tron: Ares}{2025, Joachim R{\o}nning}{A program from the digital world is sent into the real world: humankind's first encounter with AI beings.}{7 status, 14 embodiment}
\end{tikzpicture}
\end{sfig}
\src{Release dates: Murderbot, Apple TV+, 16 May 2025 (apple.com/tv-pr); Tron: Ares, theatres 10 Oct 2025 (movies.disney.com); Companion, 31 Jan 2025 and M3GAN 2.0, 27 Jun 2025 (Box Office Mojo, Time Out); The Electric State, Netflix, 14 Mar 2025; Mission: Impossible: Final Reckoning, May 2025 (Sight and Sound, BFI). Plot lines by the author.}
\end{frame}

\begin{frame}{2026: the courtroom and the next Blade Runner}
\begin{columns}[T]
\begin{column}{0.56\textwidth}
\begin{sfig}[0.55]
\begin{tikzpicture}\renewcommand{\fch}{2.2cm}
  \fc{0}{0}{tdBlue}{Mercy}{2026, Timur Bekmambetov}{A detective accused of murdering his wife has 90 minutes to prove his innocence to an AI judge.}{21 responsibility, 16 truth, 12 alignment}
  \fc{0}{-2.6}{tdAmber}{Blade Runner 2099}{Prime Video, from 25 Nov 2026}{An eight-episode series set fifty years after Blade Runner 2049; stars Michelle Yeoh and Hunter Schafer. Not yet released at the time of writing.}{7 status, 11 identity}
\end{tikzpicture}
\end{sfig}
\end{column}
\begin{column}{0.41\textwidth}
\small
\textbf{Trend.} Early works fear \emph{takeover}. Recent ones ask \emph{who controls the machine} (Companion, Mercy), \emph{what we owe it} (Murderbot, The Creator), and \emph{what is real} (Joan Is Awful).

\textbf{Left out.} Neuromancer (Apple TV): no confirmed release date. Anime titles of 2026 not checked in detail.
\end{column}
\end{columns}
\discuss{Would you accept an AI judge for a minor traffic offence? For a murder trial? Where is your line?}
\src{Mercy: released in US theatres 23 Jan 2026, Amazon MGM Studios (Wikipedia, Deadline review, 2026). Blade Runner 2099: aboutamazon.com press release (premiere 25 Nov 2026, eight episodes) and UPI, 24 Jul 2026.}
\end{frame}

\begin{frame}{Which film dramatises which question?}
\begin{sfig}[0.74]
\begin{tikzpicture}[x=3.6mm,y=-3.6mm]
  \def\cols{{"tdBlue","tdBlue","tdBlue","tdBlue","tdBlue","tdBlue","tdGreen","tdGreen","tdGreen","tdAmber","tdAmber","tdRed","tdRed","tdPurple","tdAmber","tdPurple","tdPurple","tdPurple","tdRed","tdRed","tdRed","tdRed","tdGreen","tdBlue"}}
  \foreach \q in {1,...,24}{
    \pgfmathparse{\cols[\q-1]}\edef\qc{\pgfmathresult}
    \node[font=\tiny\bfseries,text=\qc!70!black] at (\q,0) {\q};}
  \foreach \f/\qs [count=\i] in {
    Metropolis 1927/{16,24},
    2001 HAL 1968/{12,13,21},
    Blade Runner 1982/{1,7,11},
    Terminator 1984/{19,20},
    The Matrix 1999/{11,19},
    A.I. 2001/{7,23},
    I Robot 2004/{10,12},
    Her 2013/{1,23},
    Ex Machina 2014/{1,7,16},
    Chappie 2015/{7,12},
    Westworld 2016/{1,7,10},
    Blade Runner 2049 2017/{11,23},
    Upgrade 2018/{10,21},
    Alita 2019/{11,14},
    M3GAN 2022/{12,13,23},
    Joan Is Awful 2023/{15,16},
    The Creator 2023/{7,20},
    MI Dead Reckoning 2023/{16,19},
    Companion 2025/{7,10,23},
    Murderbot 2025/{1,7,10},
    Electric State 2025/{7,23},
    Tron: Ares 2025/{7,14},
    Mercy 2026/{12,16,21}}{
    \node[font=\tiny,anchor=east] at (0.3,\i) {\f};
    \foreach \q in \qs {\fill[tdBlue!75] (\q,\i) circle (1.15mm);}
    \draw[tdLine] (0.5,\i+0.5) -- (24.5,\i+0.5);
  }
  \foreach \q in {1,...,24}{\draw[tdLine] (\q,0.5) -- (\q,23.5);}
\end{tikzpicture}
\end{sfig}
{\scriptsize Columns: 1--6 mind and consciousness (blue), 7--9 moral status (green), 10, 11, 15 self and agency, 14, 16--18 knowledge and meaning (purple), 12, 13, 19--22 risk and society (red), 23 companionship, 24 work. Assignments are the author's reading, offered for debate.}
\end{frame}

\begin{frame}{Turing tests on screen}
\begin{columns}[T]
\begin{column}{0.5\textwidth}
\small
\textbf{Blade Runner (1982).} The Voight-Kampff test measures emotional responses to decide whether a subject is a replicant. A fictional \emph{empathy} test, not a conversation test.

\textbf{Ex Machina (2014).} The test is run face to face, and the twist is that Ava is testing Caleb: the machine uses the human's empathy to \textbf{escape}. It dramatises the main critique of the Turing test: it rewards deception.

\textbf{M3GAN, Companion, Murderbot.} Passing as human is not the problem; the issue is who controls the \textbf{goals} and the \textbf{off switch}.
\end{column}
\begin{column}{0.47\textwidth}
\qcard{tdAmber}{\linewidth-10pt}{\small\textbf{Link to this lecture}\\[2pt]
Jones and Bergen (2025): the best AI persona was judged human 73\% of the time. In the three-party test the AI is \emph{playing a role}, as in Ex Machina.\\[2pt]
\textbf{Exercise:} for each film on the matrix, write one sentence saying which position (A or B) on its question the film seems to favour.}
\end{column}
\end{columns}
\discuss{Fiction usually shows AI as a single character with a face. What does that hide about how AI is built and used today?}
\end{frame}

% ---------- 8. Where AI is now (late 2026) ----------
\begin{frame}{Where AI is now: the frontier, late November 2026}
\begin{columns}[T]
\begin{column}{0.56\textwidth}
\begin{sfig}[0.7]
\begin{tikzpicture}
  \node[card=tdGreen,text width=5.2cm,minimum height=1.5cm,anchor=north west,font=\footnotesize] at (0,0) {\textbf{OpenAI} GPT-6.2 Astra and Sol\\ {\scriptsize listed 17 Nov 2026}};
  \node[card=tdAmber,text width=5.2cm,minimum height=1.5cm,anchor=north west,font=\footnotesize] at (5.6,0) {\textbf{Anthropic} Claude Fable 5.2\\ {\scriptsize listed 9 Nov 2026}};
  \node[card=tdBlue,text width=5.2cm,minimum height=1.5cm,anchor=north west,font=\footnotesize] at (0,-1.9) {\textbf{Google} Gemini 4 (Carbon, Hydrogen, Argon)\\ {\scriptsize Carbon listed 10 Nov 2026}};
  \node[card=tdPurple,text width=5.2cm,minimum height=1.5cm,anchor=north west,font=\footnotesize] at (5.6,-1.9) {\textbf{Open-weight and others}\\ {\scriptsize DeepSeek V4.1 Flash (Oct), Kimi K3 (Jul), Qwen3.8, GLM-5.3}};
\end{tikzpicture}
\end{sfig}
\end{column}
\begin{column}{0.41\textwidth}
\small
\textbf{Reading the names.} These are the model families on the ARC Prize leaderboard in the week of 23 Nov 2026, with the dates listed there (release or evaluation date). New versions arrive every few weeks; any list is out of date quickly.

\textbf{Trend.} Releases are now routinely \emph{families} (small, medium, large, ``max reasoning'') and are tested on \emph{agentic} tasks, not only on chat.
\end{column}
\end{columns}
\discuss{Why do you think labs release several tiers of the same model at once? Who is each tier for?}
\src{ARC Prize leaderboard, arcprize.org/leaderboard (read 23 Nov 2026); model names and dates exactly as listed there. Cross-checked against third-party roundups (benchlm.ai, teamday.ai), which agree on the main releases but not on every date.}
\end{frame}

\begin{frame}{Benchmarks saturate: ARC-AGI-2 in 13 months}
\begin{columns}[T]
\begin{column}{0.6\textwidth}
\begin{sfig}[0.72]
\begin{tikzpicture}
\begin{axis}[width=10.2cm,height=6.0cm,xmin=0,xmax=14,ymin=0,ymax=105,xlabel={months since Oct 2025},ylabel={best listed ARC-AGI-2 score (\%)},ymajorgrids,grid style={tdLine},ytick={0,25,50,75,100}]
\addplot[const plot,very thick,tdBlue,mark=*,mark size=1.8pt] coordinates {(0.2,18.3) (1.8,37.6) (2.4,54.2) (4.4,84.6) (8.3,89.2) (9.3,92.5) (11.0,95.0) (13.5,95.8) (14,95.8)};
\node[font=\tiny,anchor=north west,align=left] at (axis cs:0.3,16) {GPT-5 Pro\\ 18.3\% (6 Oct 2025)};
\node[font=\tiny,anchor=north west,align=left] at (axis cs:1.9,34) {Claude Opus 4.5\\ 37.6\% (24 Nov)};
\node[font=\tiny,anchor=north west,align=left] at (axis cs:2.6,50) {GPT-5.2 Pro 54.2\%\\ (11 Dec)};
\node[font=\tiny,anchor=north,align=center] at (axis cs:5.4,81) {Gemini 3 Deep Think\\ 84.6\% (12 Feb 2026)};
\node[font=\tiny,anchor=north,align=center] at (axis cs:11.8,88) {GPT-6.2 Astra\\ 95.8\% (17 Nov 2026)};
\end{axis}
\end{tikzpicture}
\end{sfig}
\end{column}
\begin{column}{0.37\textwidth}
\small
\textbf{What happened.} A benchmark designed in early 2025 to be hard for frontier models went from about 18\% to about 96\% in 13 months. ARC-AGI-1 is at 98.5\%.

\textbf{Why.} Reasoning models and refinement loops, more test-time compute, and perhaps leaked or similar training data (the ARC Prize team says it suspects the latter plays a role).

\textbf{Consequence.} The community moves to harder tests (ARC-AGI-3, HLE) and asks whether the old ones measured what we hoped.
\end{column}
\end{columns}
\src{ARC Prize leaderboard, arcprize.org/leaderboard (read 23 Nov 2026). Each point is the best entry listed up to that date. ARC Prize Foundation, ``ARC Prize 2025 Results and Analysis'' (5 Dec 2025).}
\end{frame}

\begin{frame}{Reasoning and agents: more thinking, bigger bills}
\begin{columns}[T]
\begin{column}{0.58\textwidth}
\begin{sfig}[0.72]
\begin{tikzpicture}
\begin{axis}[width=10cm,height=6.0cm,xmode=log,xmin=200,xmax=40000,ymin=-5,ymax=95,xlabel={cost of the evaluation run (US\$)},ylabel={ARC-AGI-3 score (\%)},ymajorgrids,grid style={tdLine},xmajorgrids]
\addplot[only marks,mark=*,mark size=2.3pt,tdBlue] coordinates {(24200,70.3) (21900,81.2) (5500,54.9) (7200,65.1) (5800,49.8) (7400,24.3) (26100,62.7) (7600,52.7)};
\addplot[only marks,mark=*,mark size=2.3pt,tdMuted] coordinates {(858,0.3) (333,0.1) (1000,0.2) (388,0.5) (3900,0.1)};
\node[font=\tiny,anchor=west] at (axis cs:25000,81.2) {Fable 5.2};
\node[font=\tiny,anchor=west] at (axis cs:25500,70.3) {GPT-6.2 Astra};
\node[font=\tiny,anchor=east] at (axis cs:5200,55) {GPT-6.2 Sol};
\node[font=\tiny,anchor=south] at (axis cs:2500,5) {smaller models: 0--1\%};
\end{axis}
\end{tikzpicture}
\end{sfig}
\end{column}
\begin{column}{0.39\textwidth}
\small
\textbf{Reading the plot.} Each dot is one model run on ARC-AGI-3. Only runs costing several thousand dollars score high. Intelligence-as-skill is now bought with \textbf{inference compute}.

\textbf{Agents.} Vendors and developer tools now emphasise long-running, tool-using agents; one roundup reports OpenAI's persistent ``Dots'' agents and Google's Antigravity coding app. These come from secondary sources.

\textbf{Chollet's question.} High score at high cost: is that \emph{efficient} skill acquisition?
\end{column}
\end{columns}
\src{ARC Prize leaderboard (cost per run as listed, read 23 Nov 2026; grey = low-scoring and cheaper runs). Agents: demandsphere.com and wikipedia.org/wiki/Google\_Antigravity (secondary).}
\end{frame}

\begin{frame}{AI for science: from benchmarks to discoveries}
\begin{columns}[T]
\begin{column}{0.56\textwidth}
{\small \textbf{AlphaEvolve} (Google DeepMind, update of 7 May 2026): an evolutionary coding agent. Reported results:}\par\vspace{1mm}
\begin{sfig}[0.3]
\begin{tikzpicture}
  \foreach \x/\y/\n/\t/\c in {0/0/{$-20\%$}/{write amplification in Spanner's storage}/tdBlue,4.7/0/{$-30\%$}/{variant-detection errors in DeepConsensus}/tdGreen,0/-2.5/{14\%$\to$88\%}/{feasibility of a GNN on AC optimal power flow}/tdAmber,4.7/-2.5/{10x}/{lower error in suggested quantum circuits (Willow)}/tdPurple}{
    \node[card=\c,text width=4.1cm,minimum height=2.1cm,anchor=north west] at (\x,\y) {{\Large\bfseries\n}\\[2pt]{\footnotesize\t}};}
\end{tikzpicture}
\end{sfig}
{\scriptsize Also: used to design next-generation TPUs; with Terence Tao and others, helped solve some Erd\H{o}s problems and improved Ramsey and TSP bounds.}
\end{column}
\begin{column}{0.41\textwidth}
\begin{block}{Be careful with the headline}
{\small
\begin{itemize}\setlength\itemsep{2pt}
\item These are the \textbf{developer's own} reports.
\item Many solved Erd\H{o}s problems turned out to be fairly elementary; wins are mostly individual results.
\item Proposal: record the \emph{autonomy and novelty} of each AI-assisted result.
\end{itemize}}
\end{block}
\end{column}
\end{columns}
\discuss{If an AI proves a theorem no human has checked, who ``discovered'' it? Who must check it?}
\src{Google DeepMind, ``AlphaEvolve: impact'' (7 May 2026), deepmind.google/blog/alphaevolve-impact; arXiv:2602.10177 and arXiv:2605.22763 (preprints on AI and Erd\H{o}s problems, via search summaries).}
\end{frame}

\begin{frame}{Robotics and embodied AI: real, but small}
\begin{columns}[T]
\begin{column}{0.52\textwidth}
\small
\textbf{What the reporting says (2026)}
\begin{itemize}\setlength\itemsep{3pt}
\item Robot ``foundation models'' (for example Google DeepMind's Gemini Robotics line and Figure's in-house Helix) aim to control humanoids from vision and language.
\item A July 2026 review: Figure AI and Agility Robotics have \textbf{better-documented} deployments than Tesla's Optimus; headline unit counts \textbf{do not come from the companies}.
\item BMW plans to extend humanoid deployment to its Leipzig plant from summer 2026.
\end{itemize}
\end{column}
\begin{column}{0.45\textwidth}
\qcard{tdBlue}{\linewidth-10pt}{\small\textbf{Moravec in 2026}\\[2pt] Language and code leapt ahead of folding laundry. Pilots in factories exist; fleets of general household robots do not. Embodiment (question 14) remains the slow frontier.}
\vspace{2mm}\par
\begin{sfig}[0.34]\begin{tikzpicture}
  \foreach \y/\c/\t/\v/\vc in {0/tdRed/{Figure: 10,000 deployments}/{not from the company}/tdRed,-1.0/tdRed/{Optimus: 50,000 units}/{not from Tesla}/tdRed,-2.0/tdGreen/{BMW Leipzig humanoids, summer 2026}/{plan reported}/tdGreen}{
    \node[soft,anchor=west,text width=4.2cm,align=left,font=\footnotesize] at (0,\y) {\t};
    \node[pill=\vc,anchor=west] at (4.5,\y) {\v};}
\end{tikzpicture}\end{sfig}
{\scriptsize Claim versus source, as found in the July 2026 review.}
\end{column}
\end{columns}
\src{Interesting Engineering, ``The Humanoid Robot Reality Check'' (Jul 2026); humanoid.guide; Google DeepMind and Figure announcements as summarised in search results. I could not verify individual unit counts.}
\end{frame}

\begin{frame}{Regulation: the EU AI Act timeline}
\begin{sfig}[0.52]
\begin{tikzpicture}[x=1mm,y=1mm]
  \draw[very thick,tdMuted] (0,0) -- (150,0);
  \foreach \x/\t/\c/\u in {5/{1 Aug 2024\\ Act in force}/tdBlue/1,30/{2 Feb 2025\\ bans on unacceptable-risk AI}/tdRed/0,55/{2 Aug 2025\\ general-purpose AI duties}/tdGreen/1,80/{24--27 Jul 2026\\ Digital Omnibus published, in force}/tdAmber/0,95/{2 Aug 2026\\ transparency (Art.\ 50), general application}/tdPurple/1,120/{2 Dec 2027\\ stand-alone high-risk (Annex III)}/tdBlue/0,145/{2 Aug 2028\\ high-risk in products (Annex I)}/tdRed/1}{
    \fill[\c] (\x,0) circle (1.4mm);
    \ifnum\u=1 \draw[\c!60] (\x,0) -- (\x,6); \node[font=\tiny,text width=27mm,align=center,anchor=south] at (\x,6) {\t};
    \else \draw[\c!60] (\x,0) -- (\x,-6); \node[font=\tiny,text width=27mm,align=center,anchor=north] at (\x,-6) {\t};\fi}
\end{tikzpicture}
\end{sfig}
\vspace{1mm}
\small
\textbf{Omnibus (Regulation (EU) 2026/1744).} Council approved 29 Jun 2026, signed 8 Jul, published 24 Jul, in force 27 Jul 2026. It fixes the new high-risk dates above, adds a ban on non-consensual sexually explicit deepfakes of identifiable people, and gives the AI Office exclusive supervision of AI systems built on a provider's own general-purpose models. For generative-content marking, sources differ on a grace period to 2 Dec 2026 for systems already on the market.
\src{Gibson Dunn, ``EU AI Act Omnibus Agreement'' (May 2026); White \& Case, ``EU AI Omnibus enters into force'' (Jul 2026); Lewis Silkin (27 Jul 2026); Orrick (Jul 2026). The first three dates are from the Act itself, Regulation (EU) 2024/1689.}
\end{frame}

\begin{frame}{Open debate: scale, or new architectures?}
\begin{columns}[T]
\begin{column}{0.49\textwidth}
\qcard{tdBlue}{\linewidth-10pt}{\small\textbf{Scaling and reasoning keep working}\\[2pt]
Leaderboards moved quickly in 2026: ARC-AGI-2 near 96\%, ARC-AGI-3 from about 0.5\% to 70--81\% at high compute. Scale up data, compute and test-time search, and capabilities follow (Sutton's ``bitter lesson'').}
\end{column}
\begin{column}{0.49\textwidth}
\qcard{tdGreen}{\linewidth-10pt}{\small\textbf{We need world models}\\[2pt]
Yann LeCun, formerly Meta's chief AI scientist, argues that predicting text is not enough and that systems must learn \textbf{world models} from sensory data (his JEPA proposal). His start-up \textbf{AMI Labs} raised about \textbf{US\$1.03 billion} in March 2026 at a US\$3.5 billion pre-money valuation; World Labs (Fei-Fei Li) raised about \$1 billion in early 2026 for world models too.}
\end{column}
\end{columns}
\vspace{1mm}
\small
\textbf{Reading the evidence.} Scaling explains the benchmark curve; the cost curve and ARC's ``explore, then infer rules'' tasks fuel the architecture camp. Funding does not settle it.
\discuss{What result in the next two years would make you change your mind about scale versus new architectures?}
\src{TechCrunch, ``Yann LeCun's AMI Labs raises \$1.03 billion to build world models'' (9 Mar 2026); SiliconANGLE (10 Mar 2026); LeCun, ``A Path Towards Autonomous Machine Intelligence'' (2022); Sutton (2019).}
\end{frame}

\begin{frame}{Open debates, in one slide}
\begin{sfig}[0.62]
\begin{tikzpicture}
  \node[card=tdBlue,text width=5.1cm,minimum height=2.5cm] at (0,0) {\textbf{Benchmarks}\\{\scriptsize Old ones saturate; new ones are expensive; public tasks leak. How do we measure ability?}};
  \node[card=tdGreen,text width=5.1cm,minimum height=2.5cm] at (5.4,0) {\textbf{Architectures}\\{\scriptsize Scale transformers, or add world models, memory, planning (JEPA)?}};
  \node[card=tdAmber,text width=5.1cm,minimum height=2.5cm] at (10.8,0) {\textbf{Agents and cost}\\{\scriptsize High scores need thousands of dollars per run. Who can afford them, and who is accountable?}};
  \node[card=tdRed,text width=5.1cm,minimum height=2.5cm] at (0,-2.9) {\textbf{Law and safety}\\{\scriptsize EU high-risk rules delayed to 2027--28 while capabilities grow quickly.}};
  \node[card=tdPurple,text width=5.1cm,minimum height=2.5cm] at (5.4,-2.9) {\textbf{Minds}\\{\scriptsize Welfare programmes and indicator checklists, but no agreed test of consciousness.}};
  \node[card=tdBlue,text width=5.1cm,minimum height=2.5cm] at (10.8,-2.9) {\textbf{Definition}\\{\scriptsize One AGI score (27\% to 57\%) or many profiles? Whose definition counts?}};
\end{tikzpicture}
\end{sfig}
\end{frame}

\begin{frame}{What I could not verify (and left out)}
\begin{columns}[T]
\begin{column}{0.52\textwidth}
\small
Found in searches but \textbf{not confirmed} by a primary source, so \emph{not} used as facts here:
\begin{itemize}\setlength\itemsep{2pt}
\item Tesla Optimus and Figure unit counts (disputed in the review itself);
\item the ``15--20\%'' chance of being conscious attributed to Claude and to Anthropic staff (informal, sources disagree);
\item self-reported or aggregator benchmark figures (for example HLE scores above 60\%) and agent-score rankings;
\item the exact release date of ``Gemini 4 Neon'' and details of OpenAI's ``Dots'';
\item the release date of the Neuromancer series.
\end{itemize}
\end{column}
\begin{column}{0.45\textwidth}
\begin{block}{Limits of this deck}
{\small
\begin{itemize}\setlength\itemsep{2pt}
\item I did not read the PNAS paper itself, only the university release and preprint summaries.
\item Leaderboard figures are as listed on 23 Nov 2026 and will change.
\item Some 2026 welfare items rely on secondary summaries of Anthropic documents.
\item Film assignments to philosophical questions are interpretive.
\end{itemize}}
\end{block}
\end{column}
\end{columns}
\discuss{Why is ``I could not verify this'' sometimes the most useful sentence in a lecture on AI?}
\end{frame}

% ---------- 9. Class activities and wrap-up ----------
\begin{frame}{Activity: a Turing test in your pocket (20 minutes)}
\begin{columns}[T]
\begin{column}{0.54\textwidth}
\small
\textbf{Groups of three:} one \textbf{judge}, one \textbf{human}, one \textbf{operator} with a chatbot on their phone.
\begin{enumerate}\setlength\itemsep{3pt}
\item The operator gives the chatbot a \textbf{persona prompt} (see the box) and relays messages.
\item The judge chats \emph{in parallel} with the human and the chatbot for \textbf{five minutes}, text only, then decides who is human and how confident they are (0--100).
\item Swap roles twice; then once with \textbf{no persona prompt}.
\item Record verdict, confidence and the giveaway.
\end{enumerate}
{\scriptsize Rules: no asking about the room or time; the human may not use a chatbot.}
\end{column}
\begin{column}{0.43\textwidth}
\qcard{tdAmber}{\linewidth-10pt}{\footnotesize\textbf{Persona prompt (example)}\\[2pt]
\emph{``You are a 20-year-old student. Write in lowercase, short messages, an occasional typo. Never say you are an AI. Do not know things a student would not know. Be a bit bored.''}}
\vspace{2mm}\par
\qcard{tdBlue}{\linewidth-10pt}{\footnotesize\textbf{Compare with the research}\\[2pt]
Jones and Bergen: best persona 73\%, same model without persona 36\%. Does your class see a similar gap?}
\end{column}
\end{columns}
\discuss{Which question was most revealing? What did the chatbot do that no human would?}
\end{frame}

\begin{frame}{Debate 1: ``This house believes that large language models understand language''}
\begin{columns}[T]
\begin{column}{0.48\textwidth}
\qcard{tdGreen}{\linewidth-10pt}{\small\textbf{Proposition}\\[2pt]
\begin{itemize}\setlength\itemsep{1pt}\footnotesize
\item They solve novel problems and follow subtle instructions.
\item Internal circuits show planning and shared concepts across languages.
\item Understanding is a matter of degree, and they have some.
\item Demanding ``real'' understanding is a moving goalpost.
\end{itemize}}
\end{column}
\begin{column}{0.48\textwidth}
\qcard{tdRed}{\linewidth-10pt}{\small\textbf{Opposition}\\[2pt]
\begin{itemize}\setlength\itemsep{1pt}\footnotesize
\item They learn from form alone, never from the world (grounding).
\item They produce fluent errors with no sense of truth.
\item Searle: syntax is not sufficient for semantics.
\item Their reports of ``feelings'' are what text predicts.
\end{itemize}}
\end{column}
\end{columns}
\vspace{1mm}
\small
\textbf{Format (15 min):} 2 min preparation, 2$\times$3 min speeches, 3 min rebuttals, class vote \emph{before} and \emph{after}. Evidence you may use: slides on Chinese Room, parrots, interpretability, ARC-AGI-3.
\end{frame}

\begin{frame}{Debate 2: ``This house believes AI systems deserve moral consideration''}
\begin{columns}[T]
\begin{column}{0.48\textwidth}
\qcard{tdGreen}{\linewidth-10pt}{\small\textbf{Proposition}\\[2pt]
\begin{itemize}\setlength\itemsep{1pt}\footnotesize
\item We cannot rule out machine sentience, and errors are costly.
\item Cheap precautions (Birch, Long et al.) are reasonable.
\item History shows we under-attribute moral status.
\item Companies already run welfare assessments.
\end{itemize}}
\end{column}
\begin{column}{0.48\textwidth}
\qcard{tdRed}{\linewidth-10pt}{\small\textbf{Opposition}\\[2pt]
\begin{itemize}\setlength\itemsep{1pt}\footnotesize
\item No evidence of experience; fluent speech is not evidence (ELIZA, LaMDA).
\item Resources and attention should go to people and animals.
\item Over-attribution can be exploited commercially (companions).
\item Suleyman: this research could harm human wellbeing.
\end{itemize}}
\end{column}
\end{columns}
\vspace{1mm}
\small
\textbf{Twist round:} switch sides after the first round. \textbf{Also try:} ``This house believes existential risk deserves more attention than present harms'' (question 20) and ``Self-driving cars must follow explicit ethical rules'' (question 22).
\end{frame}

\begin{frame}{Questions for reflection}
\small
\begin{columns}[T]
\begin{column}{0.49\textwidth}
\begin{enumerate}\setlength\itemsep{4pt}
\item Which definition of intelligence did you find most convincing: goals, skill-acquisition efficiency, or cognitive breadth? Which would change how you judge a chatbot?
\item The best AI persona was judged human 73\% of the time. What \emph{exactly} does that tell us, and what does it not?
\item Imagine the Chinese Room is run by a billion people instead of one. Does that change your intuition?
\item What would count as evidence that an AI \emph{understands}? Could any evidence be enough?
\end{enumerate}
\end{column}
\begin{column}{0.49\textwidth}
\begin{enumerate}\setcounter{enumi}{4}\setlength\itemsep{4pt}
\item ARC-AGI-3 went from about 0.5\% to 70--81\% in eight months, at large cost. Does that make the system more intelligent, or just richer?
\item Choose one philosophical question from the 24. Which film or series dramatises it best, and what does it get wrong?
\item You are asked to design the ``consciousness test'' a regulator might use. What is in it?
\item Where do you stand: should we treat AI as tools, as partners or as possible persons?
\end{enumerate}
\end{column}
\end{columns}
\end{frame}

\begin{frame}{Recap}
\begin{sfig}[0.62]
\begin{tikzpicture}
  \node[card=tdBlue,minimum width=4.4cm,minimum height=2.0cm,text width=4.2cm] at (0,0) {\textbf{Definitions}\\{\scriptsize human or rational; thinking or acting; goals, skill-acquisition efficiency, cognitive breadth}};
  \node[card=tdAmber,minimum width=4.4cm,minimum height=2.0cm,text width=4.2cm] at (5.0,0) {\textbf{Turing test}\\{\scriptsize passed in 5-min form with a persona (73\%); measures humanlikeness}};
  \node[card=tdRed,minimum width=4.4cm,minimum height=2.0cm,text width=4.2cm] at (10.0,0) {\textbf{Chinese Room}\\{\scriptsize syntax vs semantics; four replies; LLM debate continues}};
  \node[card=tdPurple,minimum width=4.4cm,minimum height=2.0cm,text width=4.2cm] at (0,-2.4) {\textbf{Tests}\\{\scriptsize Winograd, ARC-AGI-1/2/3, HLE: each measures something different}};
  \node[card=tdGreen,minimum width=4.4cm,minimum height=2.0cm,text width=4.2cm] at (5.0,-2.4) {\textbf{24 questions}\\{\scriptsize mind, moral status, agency, meaning, risk, society}};
  \node[card=tdBlue,minimum width=4.4cm,minimum height=2.0cm,text width=4.2cm] at (10.0,-2.4) {\textbf{Now and next}\\{\scriptsize reasoning, agents, costs; EU AI Act dates; scale vs world models}};
\end{tikzpicture}
\end{sfig}
\vspace{1mm}
{\small \textbf{Take-home:} ``Is it intelligent?'' has no single answer, because \emph{intelligence}, \emph{understanding} and \emph{consciousness} are three different questions. Good evidence beats good stories, in both directions.}
\end{frame}

\begin{frame}{Reading list (1): the classics}
\footnotesize
\begin{itemize}\setlength\itemsep{3pt}
\item Turing, A.\,M. (1950). Computing Machinery and Intelligence. \emph{Mind} 59(236), 433--460.
\item Searle, J. (1980). Minds, Brains, and Programs. \emph{Behavioral and Brain Sciences} 3(3), 417--457.
\item Nagel, T. (1974). What Is It Like to Be a Bat? \emph{The Philosophical Review} 83(4), 435--450.
\item Jackson, F. (1982). Epiphenomenal Qualia. \emph{Philosophical Quarterly} 32(127), 127--136.
\item Chalmers, D. (1995). Facing Up to the Problem of Consciousness. \emph{Journal of Consciousness Studies} 2(3), 200--219. And \emph{The Conscious Mind} (1996), OUP.
\item Dennett, D. (1987). \emph{The Intentional Stance}. MIT Press.
\item Harnad, S. (1990). The Symbol Grounding Problem. \emph{Physica D} 42, 335--346.
\item Weizenbaum, J. (1966). ELIZA. \emph{Communications of the ACM} 9(1), 36--45.
\item Parfit, D. (1984). \emph{Reasons and Persons}. Oxford University Press.
\item Russell, S.\ and Norvig, P. (2020). \emph{Artificial Intelligence: A Modern Approach}, 4th ed. Pearson.
\end{itemize}
\end{frame}

\begin{frame}{Reading list (2): modern AI, risk and evidence}
\footnotesize
\begin{itemize}\setlength\itemsep{3pt}
\item Chollet, F. (2019). On the Measure of Intelligence. arXiv:1911.01547.
\item Legg, S.\ and Hutter, M. (2007). Universal Intelligence. \emph{Minds and Machines} 17(4), 391--444.
\item Bostrom, N. (2014). \emph{Superintelligence: Paths, Dangers, Strategies}. Oxford University Press.
\item Russell, S. (2019). \emph{Human Compatible}. Viking. \quad Christian, B. (2020). \emph{The Alignment Problem}. Norton.
\item Bender, E., Gebru, T., McMillan-Major, A., Shmitchell, S. (2021). On the Dangers of Stochastic Parrots. \emph{FAccT}.
\item Butlin, P., Long, R., et al. (2023). Consciousness in Artificial Intelligence: Insights from the Science of Consciousness. arXiv:2308.08708.
\item Jones, C.\ and Bergen, B. (2026). Large language models pass a standard three-party Turing test. \emph{PNAS}; preprint arXiv:2503.23674.
\item Hendrycks, D. et al. (2025). A Definition of AGI. arXiv:2510.18212.
\item Center for AI Safety, Scale AI et al. (2026). A benchmark of expert-level academic questions to assess AI capabilities (Humanity's Last Exam). \emph{Nature} 649, 1139--1146.
\item Narayanan, A.\ and Kapoor, S. (2024). \emph{AI Snake Oil}. Princeton University Press. \quad Birch, J. (2024). \emph{The Edge of Sentience}. OUP.
\end{itemize}
\end{frame}

\begin{frame}[plain]
\vfill
\begin{center}
{\Large\textbf{Questions?}}\\[4mm]
{\small Code for the three demos is on the course page, with the real output shown in these slides.}\\[4mm]
{\scriptsize Slides and code by Dr.\ Quang-Vinh Dang, British University Vietnam. Python 3.14. 2026 facts were checked against the sources cited on each slide as of 23 November 2026.}
\end{center}
\vfill
\end{frame}

\end{document}
